% concatenated from LC4BGRVS.tex
\documentclass[12pt, a4paper,reqno]{amsart}
\usepackage[a4paper, margin=4cm]{}
\usepackage{amssymb,amsfonts, amsthm, xfrac, amscd}
\usepackage{setspace}
\usepackage{breqn}
\usepackage{times}
\usepackage{changes}
\usepackage{adjustbox}
%%%%%%%%%%%%%%%%%%%%%%%%%%%%%%%%%%%%%%%
\usepackage{color}
\usepackage{array}
\usepackage{parskip}
\usepackage{microtype}
%\usepackage{refcheck}
%\usepackage{tabu}
%%%%%%%%%%%%%%%%%%%%%%%%%%%%%%%%%%%%%%%



%\usepackage{lineno}
%\linenumbers
%\usepackage[all]{xypic} 
\usepackage{color}
%\usepackage{bbm}
\usepackage[hidelinks]{hyperref}
%\usepackage[bookmarks]{hyperref}

\addtolength{\textheight}{3mm} \addtolength{\textwidth}{20mm}
%\textheight{8.8cm}
%\textwidth{6.8cm}
\addtolength{\oddsidemargin}{-8mm}
\addtolength{\evensidemargin}{-8mm} 
\addtolength{\topmargin}{-5mm}



\renewcommand{\baselinestretch}{.2}
\allowdisplaybreaks[4] 


% Over-full v-boxes on even pages are due to the \v{c} in author's name
\vfuzz4pt % Don't report over-full v-boxes if over-edge is small

\newtheoremstyle{myremark}     {11pt}{11pt}{}{}{\bfseries}{.}{.5em}{}

% THEOREM Environments
\newtheorem{thm}{Theorem}[section]
\newtheorem{cor}[thm]{Corollary}
\newtheorem{lem}[thm]{Lemma}
\newtheorem{pro}[thm]{Proposition}
\theoremstyle{definition}
\newtheorem{defn}[thm]{Definition}
\newtheorem{exmp}[thm]{Example}
%\theoremstyle{remark}
%\newtheorem{rem}[thm]{Remark}
\theoremstyle{myremark}
\newtheorem{rem}[thm]{Remark}

%\newtheorem*{thrm}{Theorem A}  
%\newtheorem*{therm}{Theorem B}  

% MATH 
\DeclareMathOperator{\RE}{Re}
\DeclareMathOperator{\IM}{Im}
\DeclareMathOperator{\ess}{ess}
\newcommand{\eps}{\varepsilon}
\newcommand{\To}{\longrightarrow}
\newcommand{\h}{\mathcal{H}}
\newcommand{\s}{\mathcal{S}}
\newcommand{\A}{\mathcal{A}}
\newcommand{\J}{\mathcal{J}}
\newcommand{\M}{\mathcal{M}}
\newcommand{\W}{\mathcal{W}}
\newcommand{\X}{\mathcal{X}}
\newcommand{\C}{\mathbb{C}}
\newcommand{\B}{\mathbf{B}}
\newcommand{\G}{\mathbb{G}}
\newcommand{\N}{\mathbb{N}}
\newcommand{\BOP}{\mathbf{B}}
\newcommand{\BH}{\mathbf{B}(\mathcal{H})}
\newcommand{\BK}{\mathbf{B}(\mathcal{K})}
\newcommand{\K}{\mathcal{K}}
\newcommand{\R}{\mathbb{R}}
\newcommand{\Complex}{\mathbb{C}}
\newcommand{\Field}{\mathbb{F}}
\newcommand{\RPlus}{\Real^{+}}
%\newcommand{\Polar}{\mathcal{P}_{\s}}
\newcommand{\Poly}{\mathcal{P}(E)}
\newcommand{\EssD}{\mathcal{D}}
\newcommand{\Lom}{\mathcal{L}}
\newcommand{\States}{\mathcal{T}}
\newcommand{\abs}[1]{\left\vert#1\right\vert}
\newcommand{\set}[1]{\left\{#1\right\}}
\newcommand{\seq}[1]{\left&lt;#1\right&gt;}
\newcommand{\norm}[1]{\left\Vert#1\right\Vert}
\newcommand{\essnorm}[1]{\norm{#1}_{\ess}}

\DeclareMathOperator{\id}{id}
\let\ker\relax                                        
\DeclareMathOperator{\ker}{Ker}                                                 
\let\hom\relax
\DeclareMathOperator{\hom}{Hom}                                                
\DeclareMathOperator{\GL}{GL}                                                   
\DeclareMathOperator{\SL}{SL}							
\DeclareMathOperator{\dom}{dom}							
\DeclareMathOperator{\tr}{tr}
\DeclareMathOperator{\Tr}{Tr}
\newcommand{\Aut}{\operatorname{Aut}}
\newcommand{\Inn}{\operatorname{Inn}}
\newcommand{\Out}{\operatorname{Out}}
\newcommand{\Mor}{\operatorname{Mor}}
\newcommand{\emp}{\varnothing}
\newcommand{\Rep}{\operatorname{Rep}}
\newcommand{\redu}{\text{red}}
\newcommand{\kar}{\text{char}}
\newcommand{\ab}{\text{ab}}
%================================================================================================
	
\onehalfspacing
\numberwithin{equation}{section}
%\doublespacing
\begin{document}
%\hfill{\today}
%\vspace{0.3cm}		
	

	
	\title[Linear combination of bilateral gamma random variables]{Linear combination of bilateral gamma random variables: distributional theory and approximations}
	
	
	\author[Barman and  Vellaisamy]{Kalyan Barman and  Palaniappan Vellaisamy}
	\address{\hskip-\parindent
		Kalyan Barman, Department of Mathematics, NIT Warangal,
		Warangal - 506004, India.}
	
	\email{barmankalyan@nitw.ac.in}

	
	%\author[Vellaisamy]{Palaniappan Vellaisamy}
	

		
		%\author[Ichiba]{Tomoyuki Ichiba}
		
		
		\address{\hskip-\parindent
		P Vellaisamy, Department of Statistics and Applied Probability, UC Santa Barbara, Santa Barbara, CA, 93106, USA.}
	\email{pvellais@ucsb.edu}
	%\date{\textit{December 28, 2023}}
		\subjclass[2020]{62E15; 62E17; 62P05; 60E05; 60E07}
	\keywords{Bilateral gamma distribution; Infinitely divisible distributions; L\'evy measure; Stein's method; Stock models}
	
	\begin{abstract}
		In this article, we obtain the exact distribution of a linear combination of bilateral gamma (BG) random variables (r.v.s). Next, we discuss the distributional properties of the linear combination of BG r.v.s, including probability density function, cumulant generating function and characteristic function. A Stein characterization is developed, which leads us to several distributional approximation results with explicit error bounds in both Kolmogorov and Wasserstein distances. Related limit theorems are also discussed. Furthermore, we show that the associated L\'evy processes are finite-variation processes with BG distributed increments having random parameters. Finally, we apply our results in exponential stock models.
	\end{abstract}

	\maketitle
	
\section{Introduction} \label{Intro}	
\noindent
The distribution of linear combination of independent gamma random variables (r.v.s) emerge naturally in various areas of probability and statistics, see Diaconis and Perlman \cite{DP1990}. The authors mentioned that the distribution of such a combination is not expressible in a closed form and so discussed approximations and studied tail probabilities. Exploiting the relationship between  negative binomial and gamma distributions (see \cite{EZ1980} or \cite{VS2010}), Vellaisamy and Upadhye \cite{pvns2009} derived the exact distribution of a linear combination of independent gamma r.v.s with arbitrary parameters.

\noindent
The bilateral-gamma (BG) distributions form a four-parameter family with a rich distributional structure (see \cite{bilateral0} and \cite{bilateral1}) and the properties that are useful for applications. This family includes the variance gamma (VG), symmetric variance gamma (SVG), and Laplace distributions, as special cases. In addition,  gamma and normal distributions appear as limiting cases, together with several other related distributions. The BG distributions are self-decomposable, stable under convolution, and have a simple cumulant generating function (cgf). The associated L\'evy processes are finite-variation processes making infinitely many jumps at each interval with positive length, and  their increments are BG distributed. Hence, the BG process provides a flexible framework for modelling fluctuations in financial markets, see K$\ddot{u}$chler and Tappe \cite{bilateral2}. 

\noindent
In this paper, we investigate the distribution of linear combinations of BG r.v.s. This family includes the linear combinations of gamma r.v.s,  as  limiting cases,  studied by Vellaisamy and Upadhye \cite{pvns2009}. We discuss several probablistic properties  of linear combinations of BG r.v.s, and show how they are related to the other distributions studied in the literature. They possess a number  of properties which make them useful for applications. The distribution of a linear combination of BG r.v.s is infinitely divisible, self-decomposable, stable under convolution, and has probability density function (pdf), cgf and characteristic function (cf) in closed form. We  derive also a Stein characterization for these distributions, obtain several approximation results using Stein's method, and discuss  related limit theorems. For the Kolmogorov distance, the error bounds for a sequence of compound Poisson distribution that converges to the distribution of a linear combination of BG r.v.s is ponited out.  For the Wasserstein distance, an error bound on the distance between  two such linear combinations is obtained. Finally, following K$\ddot{u}$chler and Tappe \cite{bilateral2}, we show that the associated L\'evy processes are finite-variation processes having infinitely many jumps at each interval with positive length, and  their increments are BG distributed, but with random parameters. 

\noindent
The paper is organized as follows. In Section \ref{nopre}, we discuss some preliminary results and in   Section \ref{probgdistribution},  the exact distribution of the linear combination of BG r.v.s  and their distributional properties are derived. Section \ref{BSBGD} deals with the approximation of linear combination of BG r.v.s. In Section \ref{BGpro}, we discuss, following K$\ddot{\text{u}}$chler and Tappe \cite{bilateral2},  the BG process associated
with  the linear combination of BG r.v.s and apply it to the exponential stock models.


\section{Preliminary results}\label{nopre}
\noindent
In this section, we introduce the notations and preliminaries required for this article.
%\vspace{-.5cm}
\subsection{Bilateral-gamma distributions}\label{prebg}
We start with the definition and related results of BG distributions. Let $Ga(\alpha,p)$ denote the gamma distribution with parameters $\alpha$ and $p$. A BG distribution (also known as gamma difference distribution) with parameters $\alpha,\beta,p,q>0$ is defined as the distribution of $X=X_1-X_2$, where $X_1\sim Ga(\alpha,p)$ and $X_2\sim Ga(\beta,q)$ and they are independent, see \cite{gamdiff}. The cf of a BG distribution is 
\begin{align}
\phi_{bg}(z)&= \frac{1}{(1 -iz/\alpha)^p (1 +iz/\beta)^q} ,~~z\in\mathbb{R}\label{bicf0}\\
&=\exp\left(\int_{\mathbb{R}}(e^{izu}-1)\nu_{bg}(du)  \right),~~z\in\mathbb{R},\label{e1}
\end{align}
\noindent
where $\nu_{bg}$ is the L\'evy measure given by
\begin{align}\label{bglevymeasure}
\nu_{bg}(du)=\left(\frac{p  }{u}e^{-\alpha u}\mathbf{1}_{(0,\infty)}(u)-\frac{q  }{u}e^{-\beta|u|}\mathbf{1}_{(-\infty,0)}(u)\right)du.
\end{align}
\noindent
Taking Fourier transform (FT) on \eqref{bicf0}, the pdf of the $BG(\alpha,p,\beta,q)$ distribution has the integral form (see Equation 1.4 of \cite{gamdiff})
\begin{align}
h_X(x)=\frac{1}{2\pi}\displaystyle\int_{\mathbb{R}}\frac{e^{-ixz}}{(1-iz/\alpha)^{p} (1+iz/\beta)^{q}}dz.
\end{align}


\noindent
An alternative form of the pdf of the $BG(\alpha,p,\beta,q)$ distribution follows from the convolution structure (see Equation 1.5 of \cite{gamdiff})
\begin{align}\label{newref02}
h_X(x)=\frac{\alpha^{p} \beta^{q}}{\Gamma p \Gamma q} \begin{cases}
e^{\beta x}\int_{x}^{\infty}y^{p-1}(y-x)^{q-1}e^{-(\alpha+\beta)y}dy  ,&~x\in(0,\infty)\\
e^{-\alpha x}\int_{-x}^{\infty}y^{q-1}(y+x)^{p-1}e^{-(\alpha+\beta)y}dy  ,&~x\in(-\infty,0).
\end{cases}
\end{align}

\noindent
Some notable special cases and limiting distributions of the BG family are as follows
(see \cite{bilateral0}). Let $\overset{d}{=}$ denote the equality in distribution.
\begin{enumerate}
	\item[(i)] When $q=p$, then $BG(\alpha,p,\beta,p)\overset{d}{=}VG(\alpha,\beta,p)$ (variance-gamma  distribution).
	
	\item [(ii)] When $\beta=\alpha$ and $q=p$, then $BG(\alpha,p,\alpha,p)\overset{d}{=}SVG(\alpha,p)$ (symmetric variance-gamma distribution).
		
	\item [(iii)] The limiting case as $p\to \infty$,  the $SVG( \sqrt{2p}/\alpha,p)$ converges in law to a normal $\mathcal{N}(0,\alpha^2)$ distribution.
	
	\item [(iv)] The limiting case as $\beta \to \infty$, the $BG(\alpha,p,\beta,q)$ converges in law to gamma $Ga(\alpha,p)$ distribution.
\end{enumerate}
\noindent
Note that the $BG$ distributions are infinitely divisible and self-decomposable, see \cite{bilateral0}. Let now $X\sim BG(\alpha,p,\beta,q)$. Then the $k$-th cumulant (see \cite[Section 2]{bilateral0}) is given by
\begin{align}
C_{k}(X)&=(k-1)!\left(\frac{p}{\alpha^k}+(-1)^k \frac{q}{\beta^k}  \right),~k\geq 1. \label{cuforbg}
\end{align}
\noindent
In particular, the first two cumulants (that is, mean and variance) of $X$ are 
$C_1(X)=\mathbb{E}(X)= \left(\frac{p}{\alpha}-\frac{q}{\beta}\right) \text{ and } C_{2}(X)=\text{Var}(X)=\left(\frac{p}{\alpha^{2}}+\frac{q}{\beta^{2}}\right).$ For more details of BG distributions, we refer the reader to  \cite{bilateral0} and \cite{bilateral1}.


\subsection{Function spaces and probability metrics}\label{PP2:FS}
Let $\mathbb{N} = \{1,2,\ldots\} \text{ and }
\mathbb{N}_{0} = \mathbb{N} \cup \{0\}$ denote henceforth the set of positive and 
non-negative integers, respectively. 

The Schwartz space is defined as
\[
\mathcal{S}(\mathbb{R})
  := \Big\{ f \in C^{\infty}(\mathbb{R}) :
        \lim_{|x|\to\infty} |x^{m} f^{(n)}(x)| = 0,
        \ \ \forall\, m,n \in \mathbb{N}_{0}
     \Big\},
\]
where $f^{(n)}$ denotes the $n$-th derivative of $f$, with the convention $f^{(0)} = f$, and 
$C^\infty(\mathbb{R})$ is the space of smooth functions on $\mathbb{R}$. It is well known that the Forier-transform (FT) is an automorphism on $\mathcal{S}(\mathbb{R})$.  
For $f \in \mathcal{S}(\mathbb{R})$, we define its FT by $\widehat{f}(u)
   := \int_{\mathbb{R}} e^{-iux} f(x)\, dx,
    u \in \mathbb{R},$ and its inverse FT by $f(x)
   := \frac{1}{2\pi} \int_{\mathbb{R}} e^{iux} \widehat{f}(u)\, du,
   x \in \mathbb{R},$ see Stein and Shakarchi \cite{stein}.

%\vspace{-1.1cm}
\noindent
We now introduce the notion of smooth Wasserstein distance (see \cite{k0} or \cite{MS001}). For a fixed integer $r \in \mathbb{N}$, consider the function class
\begin{equation}\label{fs1}
\mathcal{W}_r := 
\Big\{ h:\mathbb{R}\to \mathbb{R} \;\Big|\; 
    h \text{ is $r$-times differentiable and } 
    \|h^{(k)}\| \leq 1,\; k=0,1,\ldots,r 
\Big\},
\end{equation}
where $\|h\| := \sup_{x \in \mathbb{R}} |h(x)|$. Then, for two random variables $Y$ and $Z$, the smooth Wasserstein distance of order $r$ is defined by
\begin{equation}\label{smwdis}
  d_r(Y,Z):=d_{\mathcal{W}_r}(Y,Z)
   := \sup_{h \in \mathcal{W}_r} 
      \left|\mathbb{E}[h(Y)] - \mathbb{E}[h(Z)]\right|.  
\end{equation}

Special choices of the function class $\mathcal{W}_r$ lead to well-known probability metrics:
\begin{align*}
 d_K(Y,Z) 
   &= \Big\{h:\mathbb{R}\to \mathbb{R} \;\big|\; 
          h=\mathbf{1}_{(-\infty, x]},~ x \in \mathbb{R}\Big\}, \\[0.3em]
 d_W(Y,Z)  &= \Big\{h:\mathbb{R}\to \mathbb{R} \;\big|\; 
          h \text{ is 1-Lipschitz, i.e., } |h(x)-h(y)| \leq |x-y| \Big\},\\[0.3em]
      d_{TV}(Y,Z)&=\Big\{h:\mathbb{R}\to \mathbb{R} \;\big|\; 
          h=\mathbf{1}_{A},~ A \in \mathfrak{B} (\mathbb{R})\Big\},
\end{align*}
where $d_K(Y,Z)$ is the Kolmogorov distance,  
$d_W(Y,Z)$ is the Wasserstein distance, and $d_{TV}(Y,Z)$ is  the total variation distance (see \cite{k0}). Since $\mathfrak{B} (\mathbb{R})$ are the Borel subsets of $\mathbb{R}$, we note that (see, \cite[Lecture 2]{Chatterjee2007})
\begin{align}\label{reltvk}
d_K(Y,Z) \leq d_{TV}(Y,Z).
\end{align}
\noindent
It is shown in Appendix~A of \cite{k0} that, for $r \geq 2$,
\begin{align}\label{ineq1}
d_{r-1}(Y,Z) &\leq 3\sqrt{2}\, \sqrt{\,d_r(Y,Z)}, \text{ and}\\
d_{1}(Y,Z) &\leq 
    \Big(3\sqrt{2}\Big)^{\sum_{i=1}^{r-1} 2^{-(i-1)}} 
    \Big(d_r(Y,Z)\Big)^{2^{-(r-1)}}.
    \label{mre1}
\end{align}
\noindent
Moreover, there exists a universal constant $c>0$ such that
\[
d_{K}(Y,Z) \;\leq\; c\, \sqrt{\,d_{1}(Y,Z)}.
\]
Combining this with \eqref{mre1}, we obtain the useful bound
\[
d_{K}(Y,Z) 
   \;\leq\; 
   c \Big(3\sqrt{2}\Big)^{\,1-2^{1-r}}
   \Big(d_r(Y,Z)\Big)^{2^{-r}}, 
   \qquad r \geq 2.
\]

\vfill
\noindent
Finally, the following order relationship between the distances holds for all $r \geq 1$ (see \cite[equation 2.16]{k0}):
\begin{equation}\label{metrico1}
d_r(Y,Z) \;\leq\; d_1(Y,Z) \;\leq\; d_W(Y,Z).
\end{equation}
%\vspace{-1cm}
\subsection{Essence of Stein's Method}
Stein's method is based on the observation that a real-valued rv $Z$ has distribution $F_Z$ if and only if there exists a suitable operator $\mathcal{A}$, called the Stein operator, such that
$$
\mathbb{E}\big[\mathcal{A}f(Z)\big] = 0,
$$
for all $f$ belonging to a ``nice'' class of functions $\mathcal{F}$. This characterization naturally leads to the Stein equation
\begin{equation}\label{normal2}
\mathcal{A}f(x) = h(x) - \mathbb{E}[h(Z)],
\end{equation}
where $h$ is a real-valued test function. Replacing $x$ by $Y$ and taking expectations on both sides of \eqref{normal2} yields
\begin{equation}\label{normal3}
\mathbb{E}[h(Y)] - \mathbb{E}[h(Z)] = \mathbb{E}\big[\mathcal{A}f(Y)\big].
\end{equation}
The identity \eqref{normal3} plays a crucial role in Stein's method. For a given test function $h$, bounding the quantity $\big|\mathbb{E}[h(Y)] - \mathbb{E}[h(Z)]\big|$ reduces to obtaining suitable estimates on the solution $f$ of the Stein equation \eqref{normal2} together with information about the distribution of $Y$. For further details and applications of Stein's method, we refer the reader to the monograph \cite{nourdin} and the references therein.

\vspace{-.25cm}

\section{Linear combinations of BG random variables}\label{probgdistribution}
\noindent
 In this section, we discuss various properties of the distribution of linear combinations of BG r.v.s. First, we obtain the exact distribution of difference of linear combinations of gamma r.v.s with arbitrary parameters, which is a generalization of Theorem 3.1 of \cite{pvns2009}. In \cite{DP1990}, the authors mentioned that the distribution of sum of independent gamma r.v.s is not expressible in a closed form and so discussed approximations and studied tail probabilities. Using the connection between the negative binomial and gamma distributions (see \cite{EZ1980} or \cite{VS2010}), Vellaisamy and Upadhye \cite{pvns2009} obtain the exact distribution of weighted sums of independent gamma random variables with arbitrary parameters. 
 \vspace{-.5cm}
\subsection{Notations}\label{nota} Let $X_1,X_2,\ldots,X_n$ be independent random variables, where $X_j \sim Ga(\alpha_j,p_j)$ and $Y_1,Y_2,\ldots,Y_n$ be independent random variables, where $Y_j \sim Ga(\beta_j,q_j)$. Let $p=\sum_{j=1}^{n}p_j$ and $q=\sum_{j=1}^{n}q_j$. For $ 1 \le j \le n$, let $w_j^{(1)},w_j^{(2)}>0$ be positive constants. 


\noindent
Define now
\begin{align}\label{defTn}
	T_n= \sum_{j=1}^{n}(w_j^{(1)}X_j-w_j^{(2)}Y_j),~n\geq 1.
\end{align} 
\noindent Note when $w_j^{(1)}=w_j^{(2)},  1 \le j \le n,$ we get a linear combination of BG r.v.s.

Next, we introduce the following notations:
\begin{align}\label{ref1}
\alpha=\max_{1\leq j \leq n} \frac{\alpha_j}{w_j^{(1)}+\alpha_j} \text{ and } \beta=\max_{1\leq j \leq n} \frac{\beta_j}{w_j^{(2)}+\beta_j},
\end{align}
and for $i, n \in \mathbb{N}$,
\begin{align*}
	c_n=\prod_{j=1}^{n}\Big(  \frac{(1-\alpha)\alpha_j}{w_j^{(1)}\alpha}  \Big)^{p_j}, \quad  d_n=\prod_{j=1}^{n}\Big(  \frac{(1-\beta)\beta_j}{w_j^{(2)}\beta}  \Big)^{q_j} \\
	a_i=\frac{1}{i}\sum_{j=1}^{n}p_j\Big(\frac{1-(1-\alpha)\alpha_j}{w_j^{(1)}\alpha} \Big)^i, \quad   b_i=\frac{1}{i}\sum_{j=1}^{n}q_j\Big(\frac{1-(1-\beta)\beta_j}{w_j^{(2)}\beta} \Big)^i. 
\end{align*}


 Let $L_n$ and $M_n$ be two discrete r.v.s with probability distributions (see \cite{pvns2009}), for 
\begin{align}\label{ref2}
P(L_n=k)=c_n\gamma_k,~~  P(M_n=k)=d_n\delta_k, ~~k \ge 1,
\end{align}
 where $\gamma_0=1$, $\gamma_k=\frac{1}{k}\sum_{i=1}^{k}ia_i\gamma_{k-i}$ and $\delta_0=1$, $\delta_k=\frac{1}{k}\sum_{i=1}^{k}ib_i\delta_{k-i}$, for $k \ge 1.$. 


 
  
\begin{thm}\label{THM1}
	Let $T_n$ be defined as in \eqref{defTn}. Then $T_n\sim BG(\frac{\alpha}{1-\alpha}, L_n+p,\frac{\beta}{1-\beta},M_n+q)$, where $\alpha, \beta$ are defined in \eqref{ref1} and the r.v.s $L_n$ and $M_n$ have distributions defined in \eqref{ref2}.
\end{thm}
\begin{proof}
	First note that 
	\begin{align}\label{gdd0}
		T_n	=\sum_{j=1}^{n}(w_j^{(1)}X_j-w_j^{(2)}Y_j)=\sum_{j=1}^{n}w_j^{(1)}X_j-\sum_{j=1}^{n}w_j^{(2)}Y_j.
	\end{align}
\vspace{-0.5cm}	

	\noindent
	 From Theorem 3.1 of \cite{pvns2009}, we get $\sum_{j=1}^{n}w_j^{(1)}X_j \sim Ga(\frac{\alpha}{1-\alpha},L_n+p)$ and $\sum_{j=1}^{n}w_j^{(2)}Y_j \sim Ga(\frac{\beta}{1-\beta}, M_n+q)$. It follows, from the definition of BG distribution,  that $T_n \sim BG (\frac{\alpha}{1-\alpha},L_n+p,\frac{\beta}{1-\beta}, M_n+q)$, which proves the result. 
\end{proof}
\noindent
Since the parameters $L_n$ and $M_n$ are random, the distribution of $T_n$ is indeed a mixture of BG distributions.

\noindent
The following corollary easily follows.
\vspace{-0.5cm}
\begin{cor}\label{LCgamma}
Let the conditions of Theorem \ref{THM1} hold and $\alpha,\beta$ be defined in \eqref{ref1}. Let $L_n$ and $M_n$ be defined in \eqref{ref2}.
\begin{itemize}
\item[(i)] If $\beta \to 1$, as $n\to \infty$, then $T_n \overset{\mathfrak{L}}{\to} X,$ where $X\sim Ga(\frac{\alpha}{1-\alpha},L_n+p).$

\item[(ii)] If $\alpha \to 1$, as $n\to \infty$, then $T_n \overset{\mathfrak{L}}{\to} -Y,$ where $Y\sim Ga(\frac{\beta}{1-\beta},M_n+q).$
\end{itemize}
\end{cor}
\begin{proof}
$(i)$ Using Theorem \ref{THM1} together with \eqref{bicf0} the cf of $T_n \sim BG (\frac{\alpha}{1-\alpha},L_n+p,\frac{\beta}{1-\beta}, M_n+q)$ can be seen as
\begin{align}\label{bircf0}
\phi_{T_n}(z)=\sum_{j=0}^{\infty}\sum_{k=0}^{\infty}\frac{P(L_n=j)P(M_n=k)}{\left(1 -
\frac{iz(1-\alpha)}{\alpha}\right)^{p+j} \left(1 +\frac{iz(1-\beta)}{\beta}\right)^{q+k}}, ~~z\in\mathbb{R}.
\end{align}
\noindent
Next taking $\beta \to 1$
in \eqref{bircf0}, we get
\begin{align}\label{bircf1}
\phi_X(z)=\sum_{j=0}^{\infty}\frac{P(L_n=j)}{\left(1 -\frac{iz(1-\alpha)}{\alpha}\right)^{p+j} }, ~~z\in\mathbb{R}.
\end{align}
 which is the cf of $X=\sum_{j=1}^{n}w_j^{(1)}X_j$, where $X\sim Ga(\frac{\alpha}{1-\alpha},L_n+p)$ (see Theorem 3.1 of \cite{pvns2009}).

\noindent
$(ii)$ Similarly, taking $\alpha \to 1$ in \eqref{bircf0}, Part (ii) follows.  This proves the result.
\end{proof}

\begin{rem}\label{rmkTn}
 When $w_j^{(1)}=w_j^{(2)}=w_j$, $j\geq 1$, Theorem \ref{THM1} yields the convolution of $n$-independent BG r.v.s with arbitrary parameters. Note Theorem \ref{THM1} gives a simple random parameter representation for the distribution of $T_n$, which may be helpful for analytical or inferential purposes. 
 \end{rem}
 \vspace{-0.5cm}
\noindent
Next, we obtain pdf, moments, cumulants and Stein characterizations of linear combinations of BG r.v.s. 
The Stein characterization will be required when we will consider approximation for the distribution of linear combination of BG r.v.s.


\subsection{Probability density function}\label{pdfTn}
Taking FT on \eqref{bircf0}, the pdf of $T_n\sim BG (\eta,L_n+p,\xi, M_n+q)$ has the integral form
\begin{align}\label{birpdf0}
h_{T_n}(x)&=\frac{1}{2\pi}\sum_{j=0}^{\infty}\sum_{k=0}^{\infty}\int_{\mathbb{R}}\frac{P(L_n=j)P(M_n=k) e^{-ixz}}{(1 -iz/\eta)^{p+j} (1 +iz/\xi)^{q+k}}~dz.
\end{align}
\noindent
where $\eta=\frac{\alpha}{1-\alpha}$ and $\xi=\frac{\beta}{1-\beta}$.


 Let
 \begin{align*}
  k_1(x):= & \int_{x}^{\infty}y^{p+j-1}(y-x)^{q+k-1}e^{-(\eta+\xi)y}dy  ,~x\in(0,\infty)\\
   k_2(x):= &\int_{-x}^{\infty}y^{q+k-1}(y+x)^{p+j-1}e^{-(\eta+\xi)y}dy  ,~x\in(-\infty,0). 
\end{align*}

Using Theorem \ref{THM1} together with equation \eqref{newref02}, an alternative form of the pdf of $T_n \sim BG (\eta,L_n+p,\xi, M_n+q)$ follows due to convolution structure:
\begin{align}\label{denlcbg0}
\nonumber	h_{T_{n}}(x)&=\sum_{j=0}^{\infty}\sum_{k=0}^{\infty}\frac{P(L_n=j)\eta^{p+j} P(M_n=k)\xi^{q+k}}{\Gamma (p+j) \Gamma (q+k)} \big( e^{\xi x}k_1(x)\mathbf{1}_{(0,\infty)}(x)\\
\nonumber&\quad\quad\quad\quad+  e^{-\eta x}k_2(x)\mathbf{1}_{(-\infty,0)}(x)\big)\\
\nonumber	&=c_nd_n\sum_{j=0}^{\infty}\gamma_j\sum_{k=0}^{\infty}\delta_k\frac{\eta^{p+j} \xi^{q+k}}{\Gamma (p+j) \Gamma (q+k)} \big( e^{\xi x}k_1(x)\mathbf{1}_{(0,\infty)}(x)\\
&\quad\quad\quad\quad+  e^{-\eta x}k_2(x)\mathbf{1}_{(-\infty,0)}(x)\big)
\end{align}
\noindent
Let $ F(a,b,x)$ denote the confluent hypergeometric function (see \cite[equation (78)]{cam2001}), given by 
\begin{align}\label{chf1}
 F(a,b,x)=\frac{1}{\Gamma a}\displaystyle\int_{0}^{\infty}e^{-xt}t^{a-1}(1+t)^{b-a-1}dt, ~\text{Re}(a),\text{Re}(b),\text{Re}(x)>0.
\end{align}
\noindent
From \eqref{denlcbg0}, consider the first integral
\begin{align}\label{chf2}
	\nonumber I_1(x)&=\int_{x}^{\infty}y^{p+j-1}(y-x)^{q+k-1}e^{-(\eta+\xi)y}dy  ,~x\in(0,\infty)\\
	\nonumber &=e^{-(\eta+\xi)x}\int_{0}^{\infty}(x+u)^{p+j-1}u^{q+k-1}e^{-(\eta+\xi)u}du\\
	&=e^{-(\eta+\xi)x}x^{p+q+j+k-1}\displaystyle\int_{0}^{\infty}e^{-(\eta+\xi)xt}t^{q+k-1}(1+t)^{p+j-1}dt .
\end{align}
\noindent
Using \eqref{chf1} in \eqref{chf2}, we write
\begin{align}\label{chf3}
I_1(x)=e^{-(\eta+\xi)x}x^{p+q+j+k-1}\Gamma(q+k)F(q+k,p+q+j+k, (\eta+\xi)x), ~x\in (0,\infty).
\end{align}
\noindent
From \eqref{denlcbg0}, consider the second integral
\begin{align}\label{chf4}
I_2(x)=\int_{-x}^{\infty}y^{q+k-1}(y+x)^{p+j-1}e^{-(\eta+\xi)y}dy  ,~x\in(-\infty,0).
\end{align}
\noindent
Following the steps similar to previous integral, we have for $~x\in (-\infty,0)$,
\begin{align}\label{chf5}
I_2(x)=e^{(\eta+\xi)x}(-x)^{p+q+j+k-1}\Gamma(q+k)F(p+j,p+q+j+k, -(\eta+\xi)x). 
\end{align}
\noindent
Using \eqref{chf3} and \eqref{chf5} in \eqref{denlcbg0}, we can write the pdf of $T_n$ as
\begin{align}
 h_{T_{n}}(x)&=c_nd_n\sum_{j=0}^{\infty}\gamma_j\sum_{k=0}^{\infty}\delta_k\frac{\eta^{p+j} \xi^{q+k}}{\Gamma (p+j) \Gamma (q+k)} K_2(x),
\end{align}
where
\begin{align}
	K_2(x)= \small{\begin{cases}
			e^{-\eta x}x^{p+q+j+k-1}\Gamma(q+k)F(q+k,p+q+j+k, (\eta+\xi)x),  ~~x\in(0,\infty)\\
			e^{\xi x}(-x)^{p+q+j+k-1}\Gamma(q+k)F(p+j,p+q+j+k, -(\eta+\xi)x),  ~~x\in(-\infty,0).
	\end{cases}}
\end{align}


\vspace{-.5cm}
\begin{rem}
(i) Scheuer \cite{sch1988} obtained the pdf of $S_n=\sum_{i=1}^{n}X_i,$ where $X_i \sim Ga (\alpha_i,p_i)$ are independent r.v.s, in terms of double series involving higher order derivatives in the inner terms. Later, Vellaisamy and Upadhye \cite[Equation 4.4]{pvns2009} obtained a  simpler expression for the pdf of $S_n$, represented as a single series involving the distribution of the random variable (r.v.) $L_n$. 

\item[(ii)] Note that, if $\xi \to \infty$ in \eqref{birpdf0}, we get 
\begin{align}\label{denlcg}
\nonumber h_{S_n}(x)&=\frac{1}{2\pi}\sum_{j=0}^{\infty}\int_{\mathbb{R}}\frac{P(L_n=j) e^{-ixz}}{(1 -iz/\eta)^{p+j} }~dz\\
\nonumber &=\frac{1}{2\pi}\sum_{j=0}^{\infty}P(L_n=j)\eta^{p+j}\int_{\mathbb{R}}(\eta-iz)^{-(p+j)} e^{-ixz}~dz\\
\nonumber &=\sum_{j=0}^{\infty}P(L_n=j)\frac{\eta^{p+j} x^{p+j-1}e^{-\eta x} }{\Gamma (p+j)}\\
&=\sum_{j=0}^{\infty}c_n\gamma_j\frac{\eta^{p+j} x^{p+j-1}e^{-\eta x} }{\Gamma (p+j)},~x>0,
\end{align}
\noindent
which is the pdf of $S_n$, see \cite[Equation 4.4]{pvns2009}.
\item[(iii)] Similarly, if $\eta \to \infty$ in \eqref{birpdf0}, we get 
\begin{align*}
h_{S_n}(x)
&=\frac{1}{2\pi}\sum_{k=0}^{\infty}P(M_n=k)\xi^{p+j}\int_{\mathbb{R}}(\xi+iz)^{-(q+j)} e^{-ixz}~dz\\
&=\sum_{j=0}^{\infty}P(M_n=j)\frac{\xi^{q+j} (-x)^{q+j-1}e^{\xi x} }{\Gamma (q+j)}\\
&=\sum_{j=0}^{\infty}d_n\delta_j\frac{\eta^{q+j} (-x)^{q+j-1}e^{\xi x} }{\Gamma (q+j)},~x<0,
\end{align*}
\noindent
which is the pdf of $V_n=-\sum_{j=1}^{n}w_j^{(2)}Y_j$.

\end{rem}
\vspace{-.4cm}
\subsection{Moment generating function}\label{cflcbg} By Theorem \ref{THM1}, the mgf of $T_n$ is given by
\begin{align}
\label{mggd1}	\mathbb{E}(e^{zT_n}) &=\bigg(\sum_{j=0}^{\infty}P(L_n=j) \left( \frac{\eta}{\eta -z} \right)^{p+j}\bigg)\bigg(\sum_{k=0}^{\infty}P(M_n=k)\left( \frac{\xi}{\xi +z} \right)^{q+k}\bigg),~z\in\mathbb{R}\\
&=\mathbb{E}_{L_n}\bigg(\left( \frac{\eta}{\eta -z} \right)^{p+L_n}\bigg) \mathbb{E}_{M_n}\bigg(\left( \frac{\xi}{\xi +z} \right)^{q+M_n}\bigg),~z\in\mathbb{R},\label{ggd1}
\end{align}
\noindent
where $\mathbb{E}_{L_n}$ and $\mathbb{E}_{M_n}$ are expectations with respect to r.v.s $L_n$ and $M_n$, respectively. 
\vspace{-.2cm}
\subsection{Moments}
 Note from Theorem \ref{THM1} that $T_n:=S_n-V_n$, where  $S_n \sim Ga(\eta,L_n+p)$ and $V_n \sim Ga(\xi, M_n+q)$ (see \eqref{gdd0}). Hence, we have for $k\in \mathbb{N},$
\begin{align}\label{momentformula}
\nonumber	\mathbb{E}\left( T_n^k  \right)&=\mathbb{E}\left( (S_n-V_n)^k  \right)\\
\nonumber	&=\sum_{j=0}^{k}\binom{k}{j}(-1)^{j}\mathbb{E}(S_n^{k-j})\mathbb{E}(V_n^j)\\
\nonumber	&=\sum_{j=0}^{k}\binom{k}{j}(-1)^{j}\frac{1}{\eta^{k-j} \xi^j}\bigg(\sum_{l=0}^{\infty}P(L_n=l) \frac{\Gamma (l+p+k-j)}{\Gamma(l+p)}   \bigg)\\
\nonumber	&\quad\quad\quad\times\bigg(\sum_{m=0}^{\infty}P(M_n=m) \frac{\Gamma (m+q+j)}{\Gamma (m+q)}   \bigg)\\
\nonumber	&=\sum_{j=0}^{k}\binom{k}{j}(-1)^{j}\frac{1}{\eta^{k-j} \xi^j}\bigg(\sum_{l=0}^{\infty}c_n\gamma_l \frac{\Gamma (l+p+k-j)}{\Gamma(l+p)}   \bigg)\\
	&\quad\quad\quad\quad\quad\times\bigg(\sum_{m=0}^{\infty}d_n\delta_m \frac{\Gamma (m+q+j)}{\Gamma (m+q)}   \bigg).
\end{align}
%\vspace{-1.3cm}
\subsection{Characterizations}
 In this section, we derive a Stein characterization for $T_n$. Before stating the result, recall that the BG distributions are infinitely divisible and self-decomposable, see \cite{bilateral0}. Also, $T_n$ can be viewed as $n$-independent sums of BG r.v.s and it has BG distribution, but with random parameters. Hence, the distribution of $T_n$ is infinitely divisible and self-decomposable. Moreover, the L\'evy-Khintchine representation of the cf of $T_n$ as \eqref{e1} is given by
\begin{align}
\nonumber\phi_{T_n}(z)&=\mathbb{E}\bigg(e^{iz T_n} \bigg)=\mathbb{E}\bigg(e^{iz \left(\sum_{j=1}^{n} (w_j^{(1)}X_j-w_j^{(2)}Y_j\right) }\bigg),~z\in\mathbb{R}\\
\label{01cflcbg}&=\prod_{j=1}^{n}\left(\frac{\alpha_j}{\alpha_j-izw_{j}^{(1)}} \right)^{p_j}  \left(\frac{\beta_j}{\beta_j+izw_{j}^{(2)}}\right)^{q_j} ,~z\in\mathbb{R} \\
\label{lccf1}&=\prod_{j=1}^{n}\exp \left(\displaystyle\int_{\mathbb{R}}(e^{izu}-1) \nu_j(du)  \right),~z\in\mathbb{R} ~~(\text{use Frullani identity}),
\end{align} 
where $\nu_j(du)=\left(\frac{p_j}{u}e^{-\frac{\alpha_j u}{w_{j}^{(1)}}}\mathbf{1}_{(0,\infty)}(u)-\frac{q_j}{u}e^{-\frac{\beta_j |u|}{w_{j}^{(2)}}}\mathbf{1}_{(-\infty,0)}(u)\right)du.$ Also, the cf \eqref{lccf1} is
\begin{align}\label{lcbg2}
\phi_{T_{n}}(z)=\exp\left(\displaystyle\int_{\mathbb{R}}(e^{izu}-1) \nu_{T_n}(du)  \right), 	~z\in\mathbb{R},
\end{align}
\noindent
where $\nu_{T_n}$ is the L\'evy measure given by
\begin{align}\label{lcbg1}
\nu_{T_n}(du)=\frac{1}{u}\sum_{j=1}^{n}\left(p_je^{-\frac{\alpha_j u}{w_{j}^{(1)}}}\mathbf{1}_{(0,\infty)}(u)-q_je^{-\frac{\beta_j |u|}{w_{j}^{(2)}}}\mathbf{1}_{(-\infty,0)}(u) \right)du.
\end{align}
\noindent
When $n=1$, $w_{j}^{(1)}=w_{j}^{(2)}=1$, $\nu_{T_1}$ is the L\'evy measure of $BG(\alpha_1,p_1,\beta_1,q_1)$ distribution, see \cite{bilateral0}. Also, the $k$th order cumulant $C_k(T_n):=i^{-k}\frac{d^k}{dz^k}\log	\phi_{T_n}(z)|_{z=0}$ is
\begin{align}\label{cumu}
	C_k(T_n)=(k-1)! \int_{\mathbb{{R}}}u^{k}\nu_{T_n}(du),~k\in\mathbb{N}.
\end{align} 
%\vspace{-0.898cm}
\begin{pro}\label{th1}
	 Let $T_n$ be defined as in \eqref{defTn}. Then,
	\begin{align}\label{PP2:StenIdTn}
	\mathbb{E}\bigg(T_nf(T_n)-\displaystyle\int_{\mathbb{{R}}}uf(T_n+u)\nu_{T_n}(du) \bigg)=0,~~f\in\mathcal{S}(\mathbb{R}),
	\end{align}
	where $\nu_{T_n}$ is the L\'evy measure given in \eqref{lcbg1}.
\end{pro}

\vspace{-.6cm}
\begin{proof}
	Taking logarithms on both sides of (\ref{lcbg2}), and differentiating with respect to $z$,
	\begin{equation}\label{PP2:e4}
	\phi^{\prime}_{T_n}(z)=i\phi_{T_n}(z)\int_{\mathbb{R}}ue^{izu}\nu_{T_n}(du).
	\end{equation}
	
	\noindent
	Let $h_{T_n}(x)$ denote density of $T_n$. Then
	\begin{equation}\label{PP2:e5}
	\phi_{T_n}(z)=\displaystyle\int_{\mathbb{R}}e^{izx}h_{T_n}(x)dx~~\text{and}~~\phi^{\prime}_{T_n}(z)=i\displaystyle\int_{\mathbb{R}}xe^{izx}h_{T_n}(x)dx.
	\end{equation}
	
	\noindent
	Substituting \eqref{PP2:e5} into  \eqref{PP2:e4} and rearranging the integrals, we have
	\begin{align}	
	\displaystyle\int_{\mathbb{R}}xe^{izx}h_{T_n}(x)dx-\phi_{T_n}(z)\int_{\mathbb{R}}ue^{izu}\nu_{T_n}(du)=0\label{PP2:e6}
	\end{align}
	\noindent
	The second integral of \eqref{PP2:e6} can be written as
	\begin{align}
	\nonumber \left(\int_{\mathbb R}ue^{izu}\nu_{T_n}(du)\right)\phi_{T_n}(z)&=\int_{\mathbb R}\int_{\mathbb R}ue^{izu}e^{izx}h_{T_n}(x)dx\nu_{T_n}(du)\\
	\nonumber &=\int_{\mathbb R}\int_{\mathbb R}ue^{iz(u+x)}\nu_{T_n}(du)h_{T_n}(x)dx\\
	\nonumber &=\int_{\mathbb R}\int_{\mathbb R}ue^{izy}\nu_{T_n}(du)h_{T_n}(y-u)dy\\
	&=\int_{\mathbb R}e^{izx}\int_{\mathbb R}uh_{T_n}(x-u)\nu_{T_n}(du)dx.\label{PP2:e7}
	\end{align}
	\noindent
	Substituting \eqref{PP2:e7} into \eqref{PP2:e6}, we have
	\begin{align}
	\nonumber	0&=\displaystyle\int_{\mathbb{R}}xe^{izx}h_{T_n}(x)dx-\int_{\mathbb R}e^{izx}\int_{\mathbb R}uh_{T_n}(x-u)\nu_{T_n}(du)dx\\
	&=\displaystyle\int_{\mathbb{R}}e^{izx} \left( xh_{T_n}(x)-\int_{\mathbb R}uh_{T_n}(x-u)\nu_{T_n}(du) \right)dx. \label{PP2:e8}
	\end{align}
	\noindent
	Applying FT to \eqref{PP2:e8}, multiplying with $f\in \mathcal{S}(\mathbb{R}),$ and integrating over $\mathbb{R},$ we get
		\begin{align}\label{PP2:e9}
	\displaystyle\int_{\mathbb{R}}f(x) \left( xh_{T_n}(x)-\int_{\mathbb R}uh_{T_n}(x-u)\nu_{T_n}(du) \right)dx=0.
	\end{align}
	\noindent
	The second integral of \eqref{PP2:e9} can be seen as
	\begin{align}
	\nonumber \int_{\mathbb R}\int_{\mathbb R}uf(x)h_{T_n}(x-u)\nu_{T_n}(du)dx&=\int_{\mathbb R}\int_{\mathbb R}uf(y+u)h_{T_n}(y)\nu_{T_n}(du)dy\\ 
	&=\mathbb{E}\left(\int_{\mathbb R}f(T_n+u)u\nu_{T_n}(du)\right).\label{PP2:e10}
	\end{align}
	\noindent
	Substituting \eqref{PP2:e10} into \eqref{PP2:e9}, we have
	\begin{align*}
	\mathbb{E}\left(T_nf(T_n)-\displaystyle\int_{\mathbb{R}}f(T_n+u)u\nu_{T_n}(du) \right)=0,
	\end{align*}
	\noindent
	which proves \eqref{PP2:StenIdTn}.
	
	\noindent
	Assume conversely, \eqref{PP2:StenIdTn} holds for $\nu_{T_n}$ defined in \eqref{lcbg1}. For any $s \in \mathbb{R}$, let $f(x)=e^{isx}$, $x\in \mathbb{R}$, then (\ref{PP2:StenIdTn}) becomes
	\begin{align*}
	\mathbb{E}\left(T_ne^{isT_n}\right) &= \mathbb{E}\left(\int_{\mathbb R}e^{is(T_n+u)}u\nu_{T_n}(du)\right)\\
	&=\mathbb{E}\left(e^{isT_n}\int_{\mathbb R}e^{isu}u\nu_{T_n}(du)\right).
	\end{align*}
	\noindent
	Setting $\phi_{T_n}(s)=\mathbb{E}(e^{isX})$, then
	\begin{equation}\label{suff1}
	\phi_{T_n}^{\prime}(s)=i\phi_{T_n}(s)\int_{\mathbb R}e^{isu}u\nu_{T_n}(du).
	\end{equation}
	\noindent
	Integrating out the real and imaginary parts of (\ref{suff1}) leads, for any $z\geq 0$, to
	\begin{align*}
	\phi_{T_n}(z)&=\exp \left(i\int_{0}^{z}\int_{\mathbb R}e^{isu}u\nu_{T_n}(du)ds  \right)\\
	&=\exp\left(i\int_{\mathbb R}\int_{0}^{z}e^{isu}dsu\nu_{T_n}(du)  \right)\\
	&=\exp\left(\int_{\mathbb R}(e^{izu}-1)\nu_{T_n}(du)  \right).
	\end{align*}
	
	\noindent
	A similar computation for $z\leq0$ completes the derivation of cf of $T_n$. 
\end{proof}

\vspace*{-.7cm}
\subsection{Stein equation}\label{SEBGD}
Let $T_n$ be defined as in \eqref{defTn}. Then, from Proposition \ref{th1}, a Stein equation for $T_n$, is
\begin{equation}\label{PP2:e15}
\mathcal{A}f(x)=-xf(x)+\displaystyle\int_{\mathbb{R}}f(x+u)u\nu_{T_n}(du)= h(x)-\mathbb{E}(h(T_n)),
\end{equation}
\noindent
where $h\in \mathcal{M}$, a class of test functions. For a given $h\in \mathcal{M}$, we now use the semigroup technique to obtain $f_h$ that satisfies the integral equation \eqref{PP2:e15}. Barbour \cite{k8} developed this method for solving Stein equations, and Arras and Houdr\'e \cite{k0} extended it for infinitely divisible distributions with a finite first moment. Following Barbour's approach \cite{k8}, we choose a family of operators $(P_{t})_{t\geq0}$, for all $x\in\mathbb{R}$, as
\begin{equation}\label{PP2:e16}
P_{t}f(x):=\frac{1}{2\pi}\int_{\mathbb{R}}\widehat{f}(y)e^{iyxe^{-t} }\phi_t(y)dy, ~~f\in\mathcal{S}(\mathbb{{R}}),
\end{equation}
\noindent 
where $\hat{f}$ is the FT of $f$ and
\begin{align}\label{nPP2:a15}
\phi_t(y):=\frac{\phi_{T_n}(y)}{\phi_{T_n}(e^{-t}y)},~y\in \mathbb{{R}}.
\end{align}
\noindent
Since the distribution of $T_n$ is self-decomposable,  $\phi_{t}$ is a cf \cite[p.90]{sato} of some rv $X_{(t)}$, say. That is, for all $y\in \mathbb{R},$ and $t\geq 0,$
\begin{align}\label{PP2:a15}
\phi_t(y)=\displaystyle\int_{\mathbb{R}}e^{iy u}F_{X_{(t)}}(du),
\end{align}
\noindent
where $F_{X_{(t)}}$ is the law of $X_{(t)}$. Using \eqref{PP2:a15}, we get
\begin{align}
\nonumber  P_tf(x)&=\frac{1}{2\pi}\int_{\mathbb R}\int_{\mathbb{R}}\widehat{f}(y)e^{iy xe^{-t}}e^{iy u}F_{X_{(t)}}(du)dy\\
\nonumber  &=\frac{1}{2\pi}\int_{\mathbb R}\int_{\mathbb{R}}\widehat{f}(y)e^{iy (u+xe^{-t})}F_{X_{(t)}}(du)dy\\
&=\displaystyle\int_{\mathbb{R}}f(u+xe^{-t})F_{X_{(t)}}(du),\label{PP2:a17}
\end{align}
\noindent
where the last step follows by inverse FT (see Section \ref{PP2:FS}).
\begin{pro}\label{PP2:proSem}
	Let $(P_{t})_{t\geq0}$ be a family of operators defined in \eqref{PP2:e16}. Then
	\begin{itemize}
		\item [(i)] $(P_{t})_{t\geq 0}$ is a $\mathbb{C}_0$-semigroup on $\mathcal{S}(\mathbb{{R}})$.
		\item [(ii)] Its generator $L$ is given by
		\begin{align}
		\nonumber Lf(x)&=-xf^{\prime}(x)+\displaystyle\int_{\mathbb R}f^{\prime}(x+u)u\nu_{T_n}(du),~~f\in\mathcal{S}(\mathbb{R})\\
		\label{e7}&=\mathcal{A}f^{\prime}(x),
		\end{align}	
		
		\noindent
		where $\mathcal{A}$ is defined in \eqref{PP2:e15}.
	\end{itemize}
\end{pro}
\noindent
Following the steps similar to Proposition 3.8 and Lemma 3.10 of \cite{k1}, the proof follows.

\noindent
Next, we provide a solution of the Stein equation in \eqref{PP2:e15}.

\vfill
%\vspace{-0.3cm} 
\begin{thm}\label{thmsol}
	Let $T_n$ be defined as in \eqref{defTn} and $h\in\mathcal{W}_{r}$, defined in \eqref{fs1}. Then the function $f_{h}:\mathbb{R}\to \mathbb{R}$ defined by 
	\begin{equation}\label{PP2:SolSe}
	f_{h}(x):=-\displaystyle
	\int_{0}^{\infty}\frac{d}{dx}P_{t}h(x)dt,
	\end{equation}
	solves \eqref{PP2:e15}.	
\end{thm}


\vspace{-.3cm}
\begin{proof}
	Let
	 $$g_{h}(x)=-\displaystyle\int_{0}^{\infty}\bigg(P_{t}(h)(x)-\mathbb{E}h(T_n)  \bigg)dt .$$
	
	\noindent
	Then $g_{h}^{\prime}(x)=f_{h}(x).$ Also from \eqref{e7}, we get
	\begin{align}
	\nonumber\mathcal{A}f_{h}(x)&=-xf_{h}(x)+\int_{\mathbb{R}} f_{h}(x+u)u\nu_{T_n} (du)=Lg_{h}(x) \\
	\nonumber&=-\displaystyle\int_{0}^{\infty}LP_{t}(h)(x)dt\\
	\nonumber&=-\displaystyle\int_{0}^{\infty}\frac{d}{dt}P_{t}h(x)dt\text{ (see \cite[p.68]{nourdin})}\\
	\nonumber&=P_{0}h(x)-P_{\infty}h(x)\\
	\nonumber&=h(x)-\mathbb{E}h(T_n)~(\text{by Proposition \ref{PP2:proSem}}).
	\end{align}
	\noindent
	Hence, $f_{h}$ is the solution to \eqref{PP2:e15}. 	
\end{proof}
\vspace{-0.4cm}
\noindent
Next, we obtain some regularity conditions of $f_{h}$. 
\vspace{-0.4cm}
\begin{thm}\label{th3}
	For $h\in\mathcal{W}_{k+1}$, let $f_{h}$ be defined in \eqref{PP2:SolSe}. Then
	\begin{align}\label{PP2:pr1}
	\|f_{h}^{(k)}\|\leq \frac{1}{k+1} ,~k\geq 0,
	\end{align}	
	where $f^{(k)}$, $k\geq 1$, denotes the $k$th derivative of $f$ with $f^{(0)}=f$. Also,
	\begin{align}
		|xf_{h}^{\prime}(x)| \leq 2+\frac{1}{2}|\mathbb{E}(T_n)| \text{ and }|xf_{h}^{\prime\prime}(x)| \leq 2+\frac{1}{3}|\mathbb{E}(T_n)|.
	\end{align}
\end{thm}
\begin{proof}
	For $h\in \mathcal{W}_{k+1}$,
	\begin{align*}
	\|f_h\|&=\sup_{x\in\mathbb{R}} \left|-\displaystyle
	\int_{0}^{\infty}\frac{d}{dx}P_{t}h(x)dt\right|\\
	&=\sup_{x\in\mathbb{R}} \left| -\displaystyle
	\int_{0}^{\infty}e^{-t}\int_{\mathbb{R}}h^{(1)}(xe^{-t}+y)F_{X_{(t)}}(dy)dt \right|\\
	&\leq \|h^{(1)}\| \left|\int_{0}^{\infty}e^{-t}dt \right|
	= \|h^{(1)}\|\leq 1.
	\end{align*}
	\noindent
	Since $f_{h}$ is $k$-times differentiable, we have 
	\begin{align*}
	\|f_h^{(1)}\|&=\sup_{x\in\mathbb{R}} \left| -\displaystyle
	\int_{0}^{\infty}e^{-2t}\int_{\mathbb{R}}h^{(2)}(xe^{-t}+y)F_{X_{(t)}}(dy)dt \right|\\
	&\leq \|h^{(2)}\| \left|\int_{0}^{\infty}e^{-2t}dt \right|\\	
	&= \frac{\|h^{(2)}\|}{2} \leq \frac{1}{2}.
	\end{align*}
	\noindent
	 Indeed, it follows by induction that $\|f_{h}^{(k)}\|\leq \frac{1}{k+1},~k\geq 1.$
	 
	
	
	
	\vskip 1ex
	\noindent
	Now differentiating both sides of \eqref{PP2:e15} gives
	\begin{align}\label{sed1}
		-xf^{\prime}(x)-f(x)+\int_{\mathbb{{R}}}f^{\prime}(x+u)u\nu_{T_n}(du)=h^{\prime}(x).
	\end{align}
	\noindent
	So,
	\begin{align*}
		|xf^{\prime}(x)|& \leq \|f\| +\|h^{\prime}\|+\left|\int_{\mathbb{{R}}}f^{\prime}(x+u)u\nu_{T_n}(du)\right|\\
		&\leq \left( 2 + \|f^\prime\|\left|\int_{\mathbb{{R}}}u\nu_{T_n}(du)\right| \right)\\
		& =2+\frac{1}{2}|C_1(T_n)| ~(\text{see } \eqref{cumu})\\
		&=2+\frac{1}{2}|\mathbb{E}(T_n)|.
	\end{align*}
	\noindent
	Again differentiating both sides of \eqref{sed1} gives
	\begin{align}
			-xf^{\prime\prime}(x)-2f^{\prime}(x)+\int_{\mathbb{{R}}}f^{\prime\prime}(x+u)u\nu_{T_n}(du)=h^{\prime\prime}(x).
	\end{align}
	\noindent
	So,
	\begin{align*}
		|xf^{\prime\prime}(x)| &\leq 2\|f^{\prime}\| +\|h^{\prime\prime}\|+ \|f^{\prime\prime}\|\left|\int_{\mathbb{{R}}}u\nu_{T_n}(du)\right|\\
		& \leq 2 +\frac{1}{3}|C_1(T_n)|\\
		&=2+\frac{1}{3}|\mathbb{E}(T_n)|.
	\end{align*}
	\noindent
	This proves the result.
\end{proof}
%\vspace{-.5cm}
\section{Approximation of linear combinations of BG random variables}\label{BSBGD}
\noindent
In this section, we present bounds for the distributional approximations of linear combination of independent BG r.v.s to some probability distributions. We use Stein's method for probability approximations to derive our bounds.
%\vspace{-.6cm} 
\subsection{Approximation for sums}Our first result yields an error bound for the distributional approximation of two linear combinations of independent BG r.v.s. We refer the reader to \cite{Ransum2024}, \cite{roos2003} and \cite{PVBC} for a number of similar bounds for the total variation distance, between the distributions of two sums of non-negative integer-valued r.v.s. 
\begin{thm}\label{twosumsthm}
	For $w_j^{(1)},w_j^{(2)}, \pi_j^{(1)},\pi_j^{(2)}>0$ and $j\in \mathbb{N} \{ 0\}$, let $T_n^{(1)}=T_n= \sum_{j=1}^{n}(w_j^{(1)}X_j-w_j^{(2)}Y_j)$ and  $T_n^{(2)}= \sum_{j=1}^{n}(\pi_j^{(1)}X_j-\pi_j^{(2)}Y_j)$ be defined as in \eqref{defTn}. Then
	\begin{align}\label{twosumbd}
		d_W(T_n^{(1)},T_n^{(2)}) &\leq c_1 \sum_{j=1}^{n}\frac{p_j}{\sqrt{2\alpha_j}} \frac{ \big(w_j^{(1)} - \pi_{j}^{(1)}\big) }{\big(w_j^{(1)} + \pi_{j}^{(1)}\big)^{1/2}}+c_2\sum_{j=1}^{n}\frac{q_j}{\sqrt{2\beta_j}}\frac{\big(w_j^{(2)}-\pi_{j}^{(2)}\big)}{\big(w_j^{(2)}+\pi_{j}\big)^{1/2}},
	\end{align}
	where $c_1$ and $c_2$ are two positive constants.
\end{thm}
\vspace{-0.6cm}
 \begin{proof}
 	Let $(f,h)$ be the solution pair of \eqref{PP2:e15}. Then replacing $x$ by $T_n^{(1)} $ in the Stein equation and taking expectation, we get	
 	\begin{align}\label{ransum1}
 	\nonumber	\mathbb{E}[h(T_n^{(1)})]-\mathbb{E}[h(T_n^{(2)})]=&\mathbb{E}\left(-T_n^{(1)}f(T_n^{(1)})+\displaystyle\int_{\mathbb{R}}f(T_n^{(1)}+u)u\nu_{T_n^{(2)}}(du)  \right)\\
 \nonumber	=&\mathbb{E}\bigg[\left(-T_n^{(1)}f(T_n^{(1)})+\displaystyle\int_{\mathbb{R}}f(T_n^{(1)}+u)u\nu_{T_n^{(2)}}(du)  \right)\\
 	&\quad\quad-\left(-T_n^{(1)}f(T_n^{(1)})+\displaystyle\int_{\mathbb{R}}f(T_n^{(1)}+u)u\nu_{T_n^{(1)}}(du)  \right)\bigg],
 	\end{align}
 	\noindent
 	since $\mathbb{E}\left(-T_n^{(1)}f(T_n^{(1)})+\displaystyle\int_{\mathbb{R}}f(T_n^{(1)}+u)u\nu_{T_n^{(1)}}(du)  \right)=0, ~f\in\mathcal{S}(\mathbb{{R}}).$ So, for any $h\in \mathcal{M}_W$, \eqref{ransum1} can be written as
 	\begin{align}\label{2ransum}
 		\nonumber d_W(T_n^{(1)},T_n^{(2)}) &\leq \bigg|\mathbb{E}\displaystyle\int_{\mathbb{R}}f(T_n^{(1)}+u)u\nu_{T_n^{(2)}}(du)  -\displaystyle\int_{\mathbb{R}}f(T_n^{(1)}+u)u\nu_{T_n^{(1)}}(du)  \bigg|\\
 		\nonumber &=\bigg|\mathbb{E}\sum_{j=1}^{n}\bigg(p_j\displaystyle\int_{0}^{\infty}f(T_n^{(1)}+u)e^{-\alpha_j/\pi_{j}^{(1)} u}du\\
 		\nonumber &\quad - q_j\displaystyle\int_{0}^{\infty}f(T_n^{(1)}-u)e^{-\beta_j/\pi_{j}^{(2)} u}du\bigg)\\
 		\nonumber &\quad\quad -\sum_{j=1}^{n}\bigg(p_j\displaystyle\int_{0}^{\infty}f(T_n^{(1)}+u)e^{-\alpha_j/w_{j}^{(1)} u}du\\
 		\nonumber &\quad\quad\quad - q_j\displaystyle\int_{0}^{\infty}f(T_n^{(1)}-u)e^{-\beta_j/w_{j}^{(2)} u}du\bigg) \bigg|\\
 		\nonumber &\quad\quad\quad\quad\quad\quad\quad\quad\quad\quad\quad\quad(\text{split the measures } \nu_{T_n^{(1)}} \text{ and } \nu_{T_n^{(2)}})\\
 		\nonumber &=\bigg|\mathbb{E} \sum_{j=1}^{n}\bigg(p_j\displaystyle\int_{0}^{\infty}f(T_n^{(1)}+u)\left(e^{-\alpha_j/\pi_{j}^{(1)} u}-e^{-\alpha_j/w_{j}^{(1)} u}\right)du\\
 		\nonumber &\quad\quad\quad\quad\quad\quad-q_j\displaystyle\int_{0}^{\infty}f(T_n^{(1)}-u)\left(e^{-\beta_j/\pi_{j}^{(2)} u}-e^{-\beta_j/w_{j}^{(2)} u}\right)du\bigg) \bigg|\\
 		\nonumber &\leq\sum_{j=1}^{n}p_j \bigg|\bigg(\int_{0}^{\infty}\bigg(e^{-\alpha_j/\pi_{j}^{(1)} u}-e^{-\alpha_j/w_{j}^{(1)} u}\bigg)^2 du\bigg)^{1/2}\\
 		\nonumber&\quad\quad\quad \times
\bigg(\mathbb{E}\displaystyle\int_{0}^{\infty}f^{2}(T_n^{(1)}+u)du\bigg)^{1/2}\bigg|\\
\nonumber&\quad+\sum_{j=1}^{n}q_j \bigg|\bigg(\int_{0}^{\infty}\bigg(e^{-\beta_j/\pi_{j}^{(2)} u}-e^{-\beta_j/w_{j}^{(2)} u}\bigg)^2 du\bigg)^{1/2}\\
&\quad\quad\quad\quad\times
\bigg(\mathbb{E}\displaystyle\int_{0}^{\infty}f^{2}(T_n^{(1)}-u)du\bigg)^{1/2}\bigg|,
 \end{align}
 \noindent
 where the last inequality follows by the Cauchy-Schwartz inequality.
 
  Next, observe that, since $f\in\mathcal{S}(\mathbb{R})$, there exists two positive constants $c_1$ and $c_2$ such that $\bigg(\mathbb{E}\displaystyle\int_{0}^{\infty}f^{2}(T_n^{(1)}-u)du\bigg)^{1/2}<c_1$ and $\bigg(\mathbb{E}\displaystyle\int_{0}^{\infty}f^{2}(T_n^{(1)}-u)du\bigg)^{1/2}<c_2$, respectively. So, \eqref{2ransum} can be written as
 \begin{align*}
 d_W(T_n^{(1)},T_n^{(2)})& \leq c_1 \sum_{j=1}^{n}p_j \bigg|\bigg(\int_{0}^{\infty}\bigg(e^{-\alpha_j/\pi_{j}^{(1)} u}-e^{-\alpha_j/w_{j}^{(1)} u}\bigg)^2 du\bigg)^{\frac{1}{2}}\bigg|\\
 &\quad+c_2\sum_{j=1}^{n}q_j \bigg|\bigg(\int_{0}^{\infty}\bigg(e^{-\beta_j/\pi_{j}^{(2)} u}-e^{-\beta_j/w_{j}^{(2)} u}\bigg)^2 du\bigg)^{\frac{1}{2}}\bigg|\\
 &=c_1 \sum_{j=1}^{n}\frac{p_j}{\sqrt{\alpha_j}}\bigg|\frac{w_j^{(1)} + \pi_{j}^{(1)}}{2}-2\frac{w_j^{(1)}\pi_{j}^{(1)} }{w_j^{(1)}+\pi_{j}^{(1)}}  \bigg|^{1/2}\\
 &\quad\quad\quad + c_2\sum_{j=1}^{n}\frac{q_j}{\sqrt{\beta_j}}\bigg|\frac{w_j^{(2)}+\pi_{j}^{(2)}}{2}-2\frac{w_j^{(2)}\pi_{j}^{(2)} }{w_j^{(2)}+\pi_{j}^{(2)}}  \bigg|^{1/2}\\
 &=c_1 \sum_{j=1}^{n}\frac{p_j}{\sqrt{2\alpha_j}} \frac{ \big(w_j^{(1)} - \pi_{j}^{(1)}\big) }{\big(w_j^{(1)} + \pi_{j}^{(1)}\big)^{1/2}}+c_2\sum_{j=1}^{n}\frac{q_j}{\sqrt{2\beta_j}}\frac{\big(w_j^{(2)}-\pi_{j}^{(2)}\big)}{\big(w_j^{(2)}+\pi_{j}\big)^{1/2}}.
 \end{align*}
 This proves the result.	
 \end{proof}

\begin{cor}
	Under the assumption of Theorem \ref{twosumsthm}. Let $S_n^{(1)}= \sum_{j=1}^{n}w_j^{(1)}X_j$ and  $S_n^{(2)}= \sum_{j=1}^{n}\pi_j^{(1)}X_j$. If $\beta_j \to \infty$, as $n\to \infty$, then
	\begin{align*}
	d_W(S_n^{(1)},S_n^{(2)}) \leq c_1 \sum_{j=1}^{n}\frac{p_j}{\sqrt{2\alpha_j}} \frac{ \big(w_j^{(1)} - \pi_{j}^{(1)}\big) }{\big(w_j^{(1)} + \pi_{j}^{(1)}\big)^{1/2}},
	\end{align*}
	where $c_1$ is some positive constant.
\end{cor}
\begin{proof}
%\vspace{-0.499cm}
Note from \eqref{01cflcbg} that, if $\beta_j \to \infty$, as $n\to \infty$,
\begin{align*}
T_n^{(1)} \overset{\mathfrak{L}}{\to} S_n^{(1)} \text{ and }T_n^{(2)} \overset{\mathfrak{L}}{\to} S_n^{(2)}.
\end{align*}
Also, it follows \cite[Theorem 7.12]{Villani} that, if $T_n^{(1)} \overset{\mathfrak{L}}{\to} S_n^{(1)} \text{ and }T_n^{(2)} \overset{\mathfrak{L}}{\to} S_n^{(2)} \text{ as }\beta_j \to \infty,$
\begin{align*}
d_W(S_n^{(1)},S_n^{(2)})=\lim_{\beta_j \to \infty}d_W(T_n^{(1)},T_n^{(2)}).
\end{align*}
\noindent
Applying Theorem \ref{twosumsthm} and taking the limit as $\beta_j \to \infty$, we get from
\eqref{twosumbd}
\begin{align*}
d_W(S_n^{(1)},S_n^{(2)}) &\leq \lim_{\beta_j \to \infty} \bigg(c_1 \sum_{j=1}^{n}\frac{p_j}{\sqrt{2\alpha_j}} \frac{ \big(w_j^{(1)} - \pi_{j}^{(1)}\big) }{\big(w_j^{(1)} + \pi_{j}^{(1)}\big)^{1/2}}+c_2\sum_{j=1}^{n}\frac{q_j}{\sqrt{2\beta_j}}\frac{\big(w_j^{(2)}-\pi_{j}^{(2)}\big)}{\big(w_j^{(2)}+\pi_{j}\big)^{1/2}} \bigg) \\
 &=c_1 \sum_{j=1}^{n}\frac{p_j}{\sqrt{2\alpha_j}} \frac{ \big(w_j^{(1)} - \pi_{j}^{(1)}\big) }{\big(w_j^{(1)} + \pi_{j}^{(1)}\big)^{1/2}},
\end{align*}
\noindent
which proves the result.
\end{proof}
\begin{rem}
 Note that if $\alpha_j \to \infty,$ as $n\to \infty$, then $d_W(S_n^{(1)},S_n^{(2)}) \to 0.$
 Also, if $w_j^{(1)}=\pi_j^{(1)},$ as $n\to \infty$, then $d_W(S_n^{(1)},S_n^{(2)} ) \to 0$, as expected.
\end{rem}
\subsection{Limit of compound Poisson distributions}Next, we obtain, for the Kolmogorov distance $d_K$, the error bounds for a sequence of compound Poisson distributions that converges to the distribution of linear combination of BG r.v.s. Indeed, a more general result is shown to be true  in Theorem 1.2.18 of \cite{apple2009}, which shows an infinitely divisible distribution is a limit of compound Poisson distributions. Since $T_n$ is infinitely divisible, our result also follows from their result. However, Lemma \ref{approxthm1} provides the rate of convergence of this result. Before stating our result, let $Z_m$ be a compound Poisson r.v.s with cf (see \cite[Theorem 1.2.18]{apple2009})
	\begin{align}\label{cpdcf00}
	\phi_{m}(z):= \exp\bigg(m \left( \phi_{T_n}^{\frac{1}{m}}(z)-1  \right)   \bigg),~z\in\mathbb{{R}}, ~m\geq 1,
	\end{align}
	where $\phi_{T_n}(z)$ is given in \eqref{lcbg2}.  
\begin{lem}\label{approxthm1}
	Let $T_n$ be defined as in \eqref{defTn}. Then
	\begin{align}\label{Kbound1}
	d_{K}(Z_m,T_n) \leq c \bigg(\sum_{j=1}^{2}|C_{j}(T_n)|\bigg)^{\frac{2}{5}} \bigg(\frac{1}{m}\bigg)^{\frac{1}{5}},~m\geq 1,
	\end{align}
	where $c>0$ is some positive constant.
\end{lem} 

\vfill
\begin{proof}
	Let $b=\int_{-1}^{1}u\nu_{T_n}(du)$, where $\nu_{T_n}$ is defined in \eqref{lcbg1}. Then by Equation (2.6) of \cite{k0}, the distribution of $T_n$ is an infinitely divisible distribution with the triplet $b,0$ and $\nu_{T_n}$. Note that the distribution of $T_n$ is absolutely continuous with respect to the L\'evy measure with a bounded density and $\mathbb{E}|T_n|^2 < \infty$ (see \eqref{momentformula}). Recall from Proposition 4.11 of \cite{k0} that, if $T_n\sim ID(b,0,\nu_{T_n})$, that is, $T_n$ follows an infinitely divisible distribution with triplet $(b,0,\nu_{T_n})$ with cf $\phi_{T_n}$, and $Z_m$'s are compound Poisson r.v.s each of with cf as \eqref{cpdcf00}, then
	\begin{align}\label{kolbdd}
	d_{K}(Z_m,T_n)&\leq c\bigg( |\mathbb{E}[T_n]|+\int_{\mathbb{R}}u^2\nu_{T_n}(du)  \bigg)^{\frac{2}{p+4}}\bigg(\frac{1}{m}\bigg)^{\frac{1}{p+4}},
	\end{align}
	\noindent
	where $|\phi_{T_n}(z)| \displaystyle\int_{0}^{|z|}\frac{ds}{|\phi_{T_n}(s)|} \leq e^{c_0} |z|^p,$ $p\geq 1$.
	
	%\vskip 1ex
	\noindent
	Now observe that
	\begin{align*}
		c_0:=\sup_{s\in\mathbb{R}}\bigg|\displaystyle\int_{\mathbb{{R}}}(e^{isu} -1)\nu_{T_n}(du)  \bigg|< \infty.
	\end{align*}
	\noindent
	So,
	\begin{align*}
	|\phi_{T_n}(z)| \displaystyle\int_{0}^{|z|}\frac{ds}{|\phi_{T_n}(s)|}
	&\leq e^{c_0} \displaystyle\int_{0}^{|z|}ds\\
	&=e^{c_0}|z|.
	\end{align*}
	\noindent
	Hence by \eqref{kolbdd}, for $p=1$, we get
	\begin{align*}
	d_{K}(Z_m,T_n)&\leq c\bigg( |\mathbb{E}[T_n]|+\int_{\mathbb{R}}u^2\nu_{T_n}(du)  \bigg)^{\frac{2}{5}}\bigg(\frac{1}{m}\bigg)^{\frac{1}{5}}\\
	&= c \bigg(\sum_{j=1}^{2}|C_{j}(T_n)|\bigg)^{\frac{2}{5}} \bigg(\frac{1}{m}\bigg)^{\frac{1}{5}}, 
	\end{align*}
	\noindent
	$\text{ since }\int_{\mathbb{R}}u^2\nu_{T_n}(du)=C_2(T_n).$ This proves the result.
\end{proof}
\begin{rem}
Note that if $m \to \infty$, $d_{K}(Z_m,T_n)\to 0$, and so the distribution of $T_n$ is the limit of a compound Poisson distributions. 
\end{rem}
\vspace*{-0.7cm}
\subsection{Approximation for bilateral gamma distributions}Our next result presents an error bound for the distributional approximation of linear combination of independent BG r.v.s to a BG distribution. We refer the reader to \cite{dalygaunt2016} and \cite{leyreisw2017} for a number of similar bounds for comparisons of univariate distributions. 
%\vskip 1ex
\noindent
Let $f$ denote the solution of the Stein equation  \eqref{PP2:e15}. The proof of Theorem \ref{ApproxTHM} relies on the bounds for $|xf^{(k)}(x)|$, $k = 1, 2$, established in Theorem \ref{th3}. Define $g_n:=\prod_{j=1}^{n}\alpha_j \beta_j$,  $h_n:=g_n\sum_{j=1}^{n} (w_j^{(1)}w_j^{(2)}+|w_j^{(1)}\beta_j-w_j^{(2)}\alpha_j |  )/\alpha_j\beta_j$ and
$\kappa_n:=\frac{g_n}{g_n-h_n}.$
\vspace*{-0.4cm}
\begin{thm}\label{ApproxTHM}
Let $T_n$ be defined as in \eqref{defTn} such that $g_n> h_n$.
Also, let $Z_{bg} \sim BG(\alpha,p,\beta,q) $. Then,
\begin{align}\label{stsdbd}
\nonumber	d_{3}(T_n,Z_{bg}) & \leq \left(2+\frac{1}{3}|\mathbb{E}(T_n)|\right)\kappa_n \bigg|\sum_{j=1}^{n}\frac{w_{j}^{(1)}w_{j}^{(2)} }{\alpha_j \beta_j}- \frac{1}{\alpha \beta}\bigg|\\
\nonumber&\quad\quad+\left(2+\frac{1}{2}|\mathbb{E}(T_n)|\right)\kappa_n\bigg| \sum_{j=1}^{n}\left(\frac{w_{j}^{(1)} }{\alpha_j}-\frac{w_{j}^{(2)}}{\beta_j} \right)  -\bigg(\frac{1}{\alpha}-\frac{1}{\beta}\bigg) \bigg|\\
&\quad\quad\quad+\frac{1}{2}\kappa_n\left| \sum_{j=1}^{n}\frac{w_{j}^{(1)}w_{j}^{(2)}(p_j+q_j)}{\alpha_j \beta_j} -\frac{p+ q}{\alpha \beta} \right|+\kappa_n\left|\mathbb{E}(T_n) -\mathbb{E}(Z_{bg}) \right|. 
	\end{align}
\end{thm}
\vspace{-1.12cm}
\begin{proof}  
Let $(f,h)$ be the solution pair of \eqref{PP2:e15}. Then replacing $x$ by $Z_{bg} \sim BG(\alpha,p,\beta,q)$ in the Stein equation and taking expectation, we get
	\begin{align}
	\nonumber	\mathbb{E}(h(Z_{bg}))-\mathbb{E}(h(T_n))=&\mathbb{E}\left(-Z_{bg}f(Z_{bg})+\displaystyle\int_{\mathbb{R}}f(Z_{bg}+u)u\nu_{T_n}(du)  \right)\\
	\nonumber=&\mathbb{E} \bigg(-Z_{bg}f(Z_{bg})+\sum_{j=1}^{n}\bigg(p_j \displaystyle\int_{0}^{\infty}f(Z_{bg}+u)e^{-\alpha_j/w_{j}^{(1)} u}du\\
	\nonumber	
	&\quad -  q_j \displaystyle\int_{0}^{\infty}f(Z_{bg}-u)e^{-\beta_j/w_{j}^{(2)} u}du \bigg)\bigg)~~\text{(split the  measure $\nu_{T_n}$)}\\
	\nonumber=& \mathbb{E}\bigg( \left(  \sum_{j=1}^{n}\left(\frac{w_{j}^{(1)}p_j}{\alpha_j}-\frac{w_{j}^{(2)}q_j}{\beta_j}   \right) -Z_{bg} \right)f(Z_{bg})\bigg)\\
\nonumber	&+\sum_{j=1}^{n}\frac{w_{j}^{(1)}p_j}{\alpha_j}\mathbb{E} \left(\displaystyle\int_{0}^{\infty}f^{\prime}(Z_{bg}+u)e^{-\alpha_j/w_{j}^{(1)} u} du \right)\\
	\label{ibpf1}&\quad\quad+\sum_{j=1}^{n}\frac{w_{j}^{(2)}q_j}{\beta_j}\mathbb{E} \left(\displaystyle\int_{0}^{\infty}f^{\prime}(Z_{bg}-u)e^{-\beta_j/w_{j}^{(2)} u} du \right), 
	\end{align}
	\noindent
	where the last equality follows by the integration by parts formula. Note that
	\begin{align}\label{ibpf2}
	\nonumber&\sum_{j=1}^{n}\mathbb{E}\bigg(p_j \int_{0}^{\infty}f^{\prime}(Z_{bg}+u)e^{-\alpha_j/w_{j}^{(1)} u}du-	q_j\int_{0}^{\infty}f^{\prime}(Z_{bg}-u)e^{-\beta_j/w_{j}^{(2)} u}du\bigg)\\
	&\quad\quad\quad=\mathbb{E}\bigg(\displaystyle\int_{\mathbb R}uf^{\prime}(Z_{bg}+u)\nu_{T_n}(du)\bigg).
	\end{align}
	Using \eqref{ibpf2} in \eqref{ibpf1}, we get
	\begin{align}
	\nonumber\mathbb{E}(h(Z_{bg}))-\mathbb{E}(h(T_n))=&\mathbb{E} \bigg(\left(  \mathbb{E}(T_n) -Z_{bg} \right)f(Z_{bg})\bigg)\\
\nonumber	&+\sum_{j=1}^{n}\bigg(\left(\frac{w_{j}^{(1)} }{\alpha_j}-\frac{w_{j}^{(2)}}{\beta_j} \right)\mathbb{E} \bigg( p_j \int_{0}^{\infty}f^{\prime}(Z_{bg}+u)e^{-\alpha_j/w_{j}^{(1)} u}du\\
	\nonumber&\quad-	q_j\int_{0}^{\infty}f^{\prime}(Z_{bg}-u)e^{-\beta_j/w_{j}^{(2)} u}du \bigg)\\
	\nonumber&\quad\quad\quad+\frac{w_{j}^{(2)}p_j}{\beta_j}\mathbb{E} \left(\displaystyle\int_{0}^{\infty}f^{\prime}(Z_{bg}+u)e^{-\alpha_j/w_{j}^{(1)} u} du \right)\\
	\nonumber&\quad\quad\quad\quad\quad\quad+\frac{w_{j}^{(1)}q_j}{\alpha_j}\mathbb{E} \left(\displaystyle\int_{0}^{\infty}f^{\prime}(Z_{bg}-u)e^{-\beta_j/w_{j}^{(2)} u} du \right)\bigg)\\
	\nonumber=&\mathbb{E} \bigg(\left(  \mathbb{E}(T_n) -Z_{bg} \right)f(Z_{bg})\bigg)\\
\nonumber	&+\sum_{j=1}^{n}\left(\frac{w_{j}^{(1)} }{\alpha_j}-\frac{w_{j}^{(2)}}{\beta_j} \right)\mathbb{E} \left( \int_{\mathbb{R}} f^{\prime}(Z_{bg}+u) u\nu_{T_n}(du) \right)\\
	\nonumber&\quad\quad\quad+\sum_{j=1}^{n}\frac{w_{j}^{(2)}p_j}{\beta_j}\mathbb{E} \left(\displaystyle\int_{0}^{\infty}f^{\prime}(Z_{bg}+u)e^{-\alpha_j/w_{j}^{(1)} u} du \right)\\
	\label{ibpf3}&\quad\quad\quad\quad\quad\quad+\sum_{j=1}^{n}\frac{w_{j}^{(1)}q_j}{\alpha_j}\mathbb{E} \left(\displaystyle\int_{0}^{\infty}f^{\prime}(Z_{bg}-u)e^{-\beta_j/w_{j}^{(2)} u} du \right).
	\end{align}
	\noindent
	Applying integration by parts formula to the last two terms of \eqref{ibpf3}, we get	
	\begin{align}
	\nonumber\mathbb{E}(h(Z_{bg}))-\mathbb{E}(h(T_n))=&\mathbb{E} \bigg(\left(  \mathbb{E}(T_n) -Z_{bg} \right)f(Z_{bg})\bigg)\\
\nonumber	&\quad+\sum_{j=1}^{n}\bigg(\left(\frac{w_{j}^{(1)} }{\alpha_j}-\frac{w_{j}^{(2)}}{\beta_j} \right)\mathbb{E} \left( \int_{\mathbb{R}} f^{\prime}(Z_{bg}+u) u\nu_{T_n}(du) \right)\\
	\nonumber&\quad\quad+\frac{w_{j}^{(1)}w_{j}^{(2)}(p_j+q_j)}{\alpha_j \beta_j}\mathbb{E}\left(f^{\prime}(Z_{bg}) \right)\\
	\nonumber&\quad\quad+	\frac{w_{j}^{(1)}w_{j}^{(2)}}{\alpha_j \beta_j}\mathbb{E}\bigg(p_j \int_{0}^{\infty}f^{\prime\prime}(Z_{bg}+u)e^{-\alpha_j/w_{j}^{(1)} u}du\\
	\nonumber&\quad\quad\quad\quad-q_j\int_{0}^{\infty}f^{\prime\prime}(Z_{bg}-u)e^{-\beta_j/w_{j}^{(2)} u}du\bigg)\bigg)\\
	\nonumber =&\mathbb{E} \bigg(\left(  \mathbb{E}(T_n) -Z_{bg} \right)f(Z_{bg})\bigg)\\
	\nonumber	&\quad+\sum_{j=1}^{n}\bigg(\left(\frac{w_{j}^{(1)} }{\alpha_j}-\frac{w_{j}^{(2)}}{\beta_j} \right)\mathbb{E} \left( \int_{\mathbb{R}} f^{\prime}(Z_{bg}+u) u\nu_{T_n}(du) \right)\\
	\nonumber&\quad\quad+\frac{w_{j}^{(1)}w_{j}^{(2)}(p_j+q_j)}{\alpha_j \beta_j}\mathbb{E}\left(f^{\prime}(Z_{bg}) \right)\\
&\quad\quad+	\frac{w_{j}^{(1)}w_{j}^{(2)}}{\alpha_j \beta_j}\mathbb{E}\bigg(\displaystyle\int_{\mathbb{{R}}}f^{\prime\prime}(Z_{bg}+u)u\nu_{T_n}(du) \bigg)\bigg).\label{tt0}
	\end{align}
	\noindent
	Let $h_1,h_2 \in \mathcal{W}_3$ be two associated test functions with $f^\prime$ and $f^{\prime\prime}$ satisfy the Stein equation \eqref{PP2:e15}. Then from \eqref{tt0}, we get
	\begin{align}
	\nonumber\mathbb{E}(h(Z_{bg}))-\mathbb{E}(h(T_n)) =&\mathbb{E}\bigg(	\sum_{j=1}^{n}\frac{w_{j}^{(1)}w_{j}^{(2)}}{\alpha_j \beta_j}Z_{bg}f^{\prime\prime}(Z_{bg})+ \sum_{j=1}^{n}\bigg(\frac{w_{j}^{(1)}w_{j}^{(2)}(p_j+q_j)}{\alpha_j \beta_j}Z_{bg}\\
\nonumber	&+ \left(\frac{w_{j}^{(1)} }{\alpha_j}-\frac{w_{j}^{(2)}}{\beta_j} \right)\bigg)f^{\prime}(Z_{bg})+\left(  \mathbb{E}(T_n) -Z_{bg} \right)f(Z_{bg})  \bigg)\\
	\nonumber&\quad +\sum_{j=1}^{n}\left(\frac{w_{j}^{(1)} }{\alpha_j}-\frac{w_{j}^{(2)}}{\beta_j} \right)\bigg(\mathbb{E}(h_1(Z_{bg}))-\mathbb{E}(h_1(T_n))\bigg)\\
	&\quad\quad\quad+ \sum_{j=1}^{n}\frac{w_{j}^{(1)}w_{j}^{(2)}}{\alpha_j \beta_j} \bigg(\mathbb{E}(h_2(Z_{bg}))-\mathbb{E}(h_2(T_n))\bigg).\label{ibpf5}
	\end{align}
	\noindent
	Taking supremum over the functions $h \in \mathcal{W}_3,$ we set
\begin{align}
	\nonumber d_3(T_n,Z_{bg}) \leq ~&\mathbb{E}\bigg(	\sum_{j=1}^{n}\frac{w_{j}^{(1)}w_{j}^{(2)}}{\alpha_j \beta_j}Z_{bg}f^{\prime\prime}(Z_{bg})+ \sum_{j=1}^{n}\bigg(\frac{w_{j}^{(1)}w_{j}^{(2)}(p_j+q_j)}{\alpha_j \beta_j}Z_{bg}\\
	\nonumber	&+ \left(\frac{w_{j}^{(1)} }{\alpha_j}-\frac{w_{j}^{(2)}}{\beta_j} \right)\bigg)f^{\prime}(Z_{bg})+\left(  \mathbb{E}(T_n) -Z_{bg} \right)f(Z_{bg})  \bigg)\\
	\nonumber&\quad +\sum_{j=1}^{n}\left|\frac{w_{j}^{(1)} }{\alpha_j}-\frac{w_{j}^{(2)}}{\beta_j} \right| d_3(T_n,Z_{bg})\\
	&\quad\quad\quad+ \sum_{j=1}^{n}\frac{w_{j}^{(1)}w_{j}^{(2)}}{\alpha_j \beta_j}  d_3(T_n,Z_{bg}).\label{ibpf7}
	\end{align}
	\noindent
	That is,
	\begin{align}
	\nonumber&\left(1 - \sum_{j=1}^{n}\frac{|w_{j}^{(1)}\beta_j-w_{j}^{(2)}\alpha_j|+w_j^{(1)} w_j^{(2)}}{\alpha_j \beta_j} \right)d_{3}(T_n,Z_{bg})\\	
\nonumber	&\quad \leq\bigg| \mathbb{E}\bigg(	\sum_{j=1}^{n}\frac{w_{j}^{(1)}w_{j}^{(2)}}{\alpha_j \beta_j}Z_{bg}f^{\prime\prime}(Z_{bg})+ \sum_{j=1}^{n}\bigg(\frac{w_{j}^{(1)}w_{j}^{(2)}(p_j+q_j)}{\alpha_j \beta_j}Z_{bg}\\
	&\quad\quad+ \left(\frac{w_{j}^{(1)} }{\alpha_j}-\frac{w_{j}^{(2)}}{\beta_j} \right)\bigg)f^{\prime}(Z_{bg})+\left(  \mathbb{E}(T_n) -Z_{bg} \right)f(Z_{bg})  \bigg)\bigg|.\label{tt1}
	\end{align}
\noindent
Let $\mathcal{A}_{bg}$ denote a Stein operator of $Z_{bg}\sim BG(\alpha,p,\beta,q)$. Then
\begin{align}
	\mathbb{E}\left(\mathcal{A}_{bg}f(Z_{bg}) \right)=0,~f\in \mathcal{S}(\mathbb{{R}}),
\end{align}	
where $\mathcal{A}_{bg}f(x)=\frac{1}{\alpha \beta}xf^{\prime\prime}(x)+ \left(\frac{p+q}{\alpha\beta}+ \left(\frac{1}{\alpha}-\frac{1}{\beta} \right)x   \right)f^{\prime}(x)+\left(  \mathbb{E}(Z_{bg}) -x \right)f(x)$, see \cite[Example 5.3]{BUV1}. Then, from \eqref{tt1},
\begin{align}
\nonumber&\left(1 - \sum_{j=1}^{n}\frac{|w_{j}^{(1)}\beta_j-w_{j}^{(2)}\alpha_j|+w_j^{(1)} w_j^{(2)}}{\alpha_j \beta_j} \right)d_{3}(T_n,Z_{bg})\\	
\nonumber	&\quad \leq\bigg| \mathbb{E}\bigg(	\sum_{j=1}^{n}\frac{w_{j}^{(1)}w_{j}^{(2)}}{\alpha_j \beta_j}Z_{bg}f^{\prime\prime}(Z_{bg})+ \sum_{j=1}^{n}\bigg(\frac{w_{j}^{(1)}w_{j}^{(2)}(p_j+q_j)}{\alpha_j \beta_j}Z_{bg}\\
\nonumber	&\quad\quad+ \left(\frac{w_{j}^{(1)} }{\alpha_j}-\frac{w_{j}^{(2)}}{\beta_j} \right)\bigg)f^{\prime}(Z_{bg})+\left(  \mathbb{E}(T_n) -Z_{bg} \right)f(Z_{bg})  \bigg)- \mathbb{E}\bigg(\frac{1}{\alpha \beta}Z_{bg}f^{\prime\prime}(Z_{bg})\\
\nonumber&\quad\quad\quad\quad+ \left(\frac{p+q}{\alpha\beta}+ \left(\frac{1}{\alpha}-\frac{1}{\beta} \right)Z_{bg}   \right)f^{\prime}(Z_{bg})+\left(  \mathbb{E}(Z_{bg}) -Z_{bg} \right)f(Z_{bg})  \bigg) \bigg|\\
\nonumber&\quad=\bigg|\mathbb{E}\bigg( \left( \sum_{j=1}^{n}\frac{w_{j}^{(1)}w_{j}^{(2)} }{\alpha_j \beta_j}- \frac{1}{\alpha \beta} \right)Z_{bg}f^{\prime\prime}(Z_{bg}) \\
\nonumber &\quad+ \left(\sum_{j=1}^{n}\left(\frac{w_{j}^{(1)} }{\alpha_j}-\frac{w_{j}^{(2)}}{\beta_j} \right)  -\bigg(\frac{1}{\alpha}-\frac{1}{\beta}\bigg)\right)Z_{bg}f^{\prime}(Z_{bg}) \bigg) \\
\nonumber&\quad\quad+\left( \sum_{j=1}^{n}\frac{w_{j}^{(1)}w_{j}^{(2)}(p_j+q_j)}{\alpha_j \beta_j} -\frac{p+ q}{\alpha \beta} \right)f^{\prime}(Z_{bg})+ \left(\mathbb{E}(T_n) -\mathbb{E}(Z_{bg}) \right)f(Z_{bg}) \bigg|\\
\nonumber& \leq \bigg|\sum_{j=1}^{n}\frac{w_{j}^{(1)}w_{j}^{(2)} }{\alpha_j \beta_j}- \frac{1}{\alpha \beta}\bigg|\|xf^{\prime\prime}(x)\|+\bigg| \sum_{j=1}^{n}\left(\frac{w_{j}^{(1)} }{\alpha_j}-\frac{w_{j}^{(2)}}{\beta_j} \right)  -\bigg(\frac{1}{\alpha}-\frac{1}{\beta}\bigg) \bigg|\|xf^{\prime}(x)\|\\
\nonumber &\quad\quad\quad+\left| \sum_{j=1}^{n}\frac{w_{j}^{(1)}w_{j}^{(2)}(p_j+q_j)}{\alpha_j \beta_j} -\frac{p+ q}{\alpha \beta} \right|\|f^{\prime}\|+\left|\mathbb{E}(T_n) -\mathbb{E}(Z_{bg}) \right|\|f\|,
\end{align}
\noindent
hence,
\begin{align}\label{lastbd}
\nonumber d_3(T_n,Z_{bg})&\leq \kappa_n \bigg(\bigg|\sum_{j=1}^{n}\frac{w_{j}^{(1)}w_{j}^{(2)} }{\alpha_j \beta_j}- \frac{1}{\alpha \beta}\bigg|\|xf^{\prime\prime}(x)\|+\bigg| \sum_{j=1}^{n}\left(\frac{w_{j}^{(1)} }{\alpha_j}-\frac{w_{j}^{(2)}}{\beta_j} \right) \\
\nonumber&\quad\quad
 -\bigg(\frac{1}{\alpha}-\frac{1}{\beta}\bigg) \bigg|\|xf^{\prime}(x)\|+\left| \sum_{j=1}^{n}\frac{w_{j}^{(1)}w_{j}^{(2)}(p_j+q_j)}{\alpha_j \beta_j} -\frac{p+ q}{\alpha \beta} \right|\|f^{\prime}\|\\
 &\quad\quad\quad\quad+\left|\mathbb{E}(T_n) -\mathbb{E}(Z_{bg}) \right|\|f\|  \bigg).
\end{align}
\noindent
Using the estimates of Theorem \ref{th3} to bound \eqref{lastbd} yields \eqref{stsdbd}.
\end{proof}
\begin{rem}
 Note that if $w_j^{(1)}=w_j^{(2)}=1$, $n=1$, $\alpha_j \to \alpha,$ $\beta_j \to \beta$, $p_j \to p$ and $q_j \to q$, then by \eqref{stsdbd}, $d_3(T_1,Z_{bg}) \to 0$, as expected.
\end{rem}
\vspace{-0.5cm}
\noindent
We now have the following corollary for VG approximation to the linear combinations of independent BG r.v.s. The VG approximation for sums of independent r.v.s has been widely studied in the literature; see, for instance, \cite{Gaunt2014, Gaunt2020, Gaunt2022}.
\vspace{-0.5cm}
\begin{cor}\label{ApproxTHM1}
Let $T_n$ be defined as in \eqref{defTn} such that $g_n> h_n$.
Also let $Z_{vg} \sim VG (\alpha,\beta,p)$. Then,
\begin{align}\label{stsdbdvg}
\nonumber d_{3}(T_n,Z_{vg}) & \leq \left(2+\frac{1}{3}|\mathbb{E}(T_n)|\right)\kappa_n\bigg|\sum_{j=1}^{n}\frac{w_{j}^{(1)}w_{j}^{(2)} }{\alpha_j \beta_j}- \frac{1}{\alpha \beta}\bigg|\\
\nonumber &\quad\quad+\left(2+\frac{1}{2}|\mathbb{E}(T_n)|\right)\kappa_n\bigg| \sum_{j=1}^{n}\left(\frac{w_{j}^{(1)} }{\alpha_j}-\frac{w_{j}^{(2)}}{\beta_j} \right)  -\bigg(\frac{1}{\alpha}-\frac{1}{\beta}\bigg) \bigg|\\
&\quad\quad\quad+\frac{1}{2}\kappa_n\left| \sum_{j=1}^{n}\frac{w_{j}^{(1)}w_{j}^{(2)}(p_j+q_j)}{\alpha_j \beta_j} -\frac{2p}{\alpha \beta} \right|+\kappa_n\left|\mathbb{E}(T_n) -\mathbb{E}(Z_{vg}) \right|.
\end{align}
\end{cor}
\begin{rem}
Note that if $w_j^{(1)}=w_j^{(2)}=1$, $p_j=q_j=p_n$, $\alpha_j=\beta_j=\sqrt{n} \alpha$, $\alpha=\beta$, and $Var(T_n) \to Var (Z_{vg})$, then by \eqref{stsdbdvg}, $d_3(T_n,Z_{vg}) \to 0$ as $n \to \infty$. 
\end{rem}
\vspace{-0.5cm}
\noindent
The following example establishes an explicit error bound for the normal approximation of linear combinations of BG r.v.s.
\vspace{-0.6cm}
\begin{exmp}
	Let $Z_p\sim SVG(\sqrt{2p}/\sigma, p)$ and $Z_\sigma \sim \mathcal{N}(0,\sigma^2) $. Recall from Section \ref{prebg} that,
$SVG(\sqrt{2p}/\sigma, p) \overset{d}{=}BG(\sqrt{2p}/\sigma, p,\sqrt{2p}/\sigma, p)$. Then, the cf of $SVG(\sqrt{2p}/\sigma, p)$ is
\begin{align}
\label{svgcf000}	\phi_{sv}(z)&= \bigg(1+ \frac{z^2 \sigma^2}{2p} \bigg)^{-p},~z\in\mathbb{R}\\
&=\exp\bigg(\displaystyle\int_{\mathbb{R}} (e^{izu}-1)\nu_{sv}(du)  \bigg),~z\in\mathbb{{R}},
\end{align}  
\noindent
where the L\'evy measure $\nu_{sv}$ is 
\begin{align*}
\nu_{sv}(du)=\bigg(\frac{p}{u}e^{- \frac{\sqrt{2p}}{\sigma} u}\textbf{1}_{(0,\infty)}(u)+ \frac{p}{|u|}e^{- \frac{\sqrt{2p}}{\sigma} |u|}\textbf{1}_{(-\infty,0)}(u) \bigg).
\end{align*}
\noindent
Note from \eqref{svgcf000} that, 
\begin{align*}
\lim_{p\to\infty}\phi_{sv}(z)=e^{-\frac{\sigma^2 z^2}{2}}.
\end{align*}
\noindent
That is, $Z_p \overset{L}{\to} Z_\sigma \sim \mathcal{N}(0,\sigma^2),$ as $p\to\infty.$ Also, it follows \cite[Theorem 7.12]{Villani} that, if $Z_p \overset{L}{\to} Z_\sigma,$ as $p \to \infty$,


\begin{align}\label{svgcf1100}
d_{3}(T_n,Z_\sigma)= \lim_{p \to \infty}d_{3}(T_n,Z_p).
\end{align}
\noindent
Applying Theorem \ref{ApproxTHM} to $Z_{bg} = Z_p$, and taking the limit as $p\to \infty$, we get from \eqref{stsdbd}
	\begin{align}\label{bgclt}
\nonumber	d_{3}(T_n,Z_\sigma) & \leq \lim_{p\to\infty}\bigg(  \left(2+\frac{1}{3}|\mathbb{E}(T_n)|\right)\kappa_n \bigg|\sum_{j=1}^{n}\frac{w_{j}^{(1)}w_{j}^{(2)} }{\alpha_j \beta_j}- \frac{\sigma^2}{2p}\bigg|\\
\nonumber &\quad\quad+\left(2+\frac{1}{2}|\mathbb{E}(T_n)|\right)\kappa_n\bigg| \sum_{j=1}^{n}\left(\frac{w_{j}^{(1)} }{\alpha_j}-\frac{w_{j}^{(2)}}{\beta_j} \right)\bigg|\\
\nonumber &\quad\quad\quad+\frac{1}{2}\kappa_n\left| \sum_{j=1}^{n}\frac{w_{j}^{(1)}w_{j}^{(2)}(p_j+q_j)}{\alpha_j \beta_j} -\sigma^2 \right|+\kappa_n\left|\mathbb{E}T_n\right|\bigg)\\
\nonumber &= \left(2+\frac{1}{3}|\mathbb{E}(T_n)|\right)\kappa_n \sum_{j=1}^{n}\frac{w_{j}^{(1)}w_{j}^{(2)} }{\alpha_j \beta_j}\\
\nonumber &\quad\quad+\left(2+\frac{1}{2}|\mathbb{E}(T_n)|\right)\kappa_n\bigg| \sum_{j=1}^{n}\left(\frac{w_{j}^{(1)} }{\alpha_j}-\frac{w_{j}^{(2)}}{\beta_j} \right)\bigg|\\
&\quad\quad\quad+\frac{1}{2}\kappa_n\left| \sum_{j=1}^{n}\frac{w_{j}^{(1)}w_{j}^{(2)}(p_j+q_j)}{\alpha_j \beta_j} -\sigma^2 \right|+\kappa_n\left|\mathbb{E}T_n\right|,
\end{align}
which gives the error in the closeness between the distributions of $T_n$ and $Z_\sigma$. 
\end{exmp}
\begin{rem}
(i) Note that if we set $w_j^{(1)}=w_j^{(2)}=w_j$, $p_j=q_j$, $\alpha_j=\beta_j$ such that $w_j/\alpha_j,w_j/\beta_j \to 0$ and $Var(T_n) \to \sigma
^2$, then by \eqref{bgclt}, $d_3(T_n, Z_\sigma) \to 0,$ as $n \to \infty$.

\item [(ii)]Note also that, if $\sigma=1, w_j^{(1)}=w_j^{(2)}=1/\sqrt{n}, \alpha_j=\beta_j=n\alpha , \text{ and } p_j=q_j=\frac{1}{2}n^3\alpha^2$ in \eqref{bgclt}, we get $d_3(T_n,Z_1) \to 0$ as $n\to \infty$. This result is, in a sense a generalization of the Proposition 3.1 of \cite{bilateral0}, since $T_1=w_1^{(1)}X_1-w_1^{(2)}Y_1\sim $ BG distribution.
\end{rem}


\section{Mixed Bilateral gamma process}\label{BGpro}
\noindent
In this section, we first introduce the mixed bilateral gamma process, which is a BG process with random parameters. We then discuss some of its important properties.

\noindent
Let $\{T_n(t)\}_{t\geq 0}$ be a L\'evy process with cf 
\begin{align}\label{cfBGRP}
	\mathbb{E}(e^{izT_n(t)})=e^{t \tau (z)},~z\in\mathbb{{R}},
\end{align}
\noindent
where $\tau(z)$ is the characteristic exponent, given by $\tau(z)=\int_{\mathbb{{R}}}(e^{izu}-1)\nu_{T_n}(du)$ and $\nu_{T_n}$ is the L\'evy measure given in \eqref{lcbg1}. Note the case $n=1$ was first considered by K$\ddot{\text{u}}$chler and Tappe \cite{bilateral0}. Recently, by incorporating a diffusion component into the BG process, Kirkby {\it et al.} \cite{bgm2024} introduced BG motion and applied it in stock models. From \eqref{lcbg1}, we see that $\nu_{T_n}(\mathbb{R})=\infty$ and $\int_{\mathbb{R}}|u| \nu_{T_n}(du)<\infty$. Since the Gaussian component is zero, $T_n(t)$ is classified as type B according to Definition 11.9 of Sato \cite{sato}. Moreover, $T_n(t)$ is a finite-variation process (see \cite[Theorem 21.9]{sato}) with countable jumping times, almost surely (see \cite[Theorem 21.3]{sato}). We note that all the increments of $T_n(t)$ have the BG distribution, but with random parameters. More precisely, 
\begin{align}\label{bgprorandom}
T_n(t) -T_n(s)\sim BG\left(\eta, (L_n+p)(t-s),\xi,(M_n+q)(t-s)\right),~0\leq s \leq t,	
\end{align}
\noindent
 where $\eta=\alpha/1-\alpha$, $\xi=\beta/1-\beta$, and $\alpha,\beta,L_n,M_n,p$, and $q$ are as defined in \eqref{ref1} and \eqref{ref2}, respectively. Also,
 \begin{align}\label{lcbgpro1}
T_n(t) \sim BG\left(\eta, (L_n+p)t,\xi,(M_n+q)t\right),~ t\geq 0.	
\end{align}
\noindent
Since $L_n$ and $M_n$ are random,  $T_n(t)$ is indeed a mixed BG process.

 \noindent
 Next, we present an application of the above introduced process to stock models.
 \vspace{-0.5cm}
 \subsection{Arbitrage-free Stock models}
 In the continuous time finance, one of the crucial task is to find a realistic and analytically tractable models for price evolutions of risky financial assets (see, \cite{bgm2024} and \cite{bilateral0}, \cite{bilateral2}). We consider the exponential L\'evy models (see, K$\ddot{\text{u}}$chler and Tappe \cite{bilateral2,tsd2013})
  \begin{align}\label{LCBGstock}
 	\begin{cases}
 	S_t&=S_0e^{T_n(t)}\\
 	B_t&=e^{rt},
 	\end{cases}
 \end{align}
\noindent
which consists two financial assets $(S,B)$. Here $S$ is the dividend paying stock with deterministic initial value $S_0$ with dividend rate $v\geq 0$. Also, $B$ is the bank account with fixed interest rate $r\geq 0$.  Note that the model \eqref{LCBGstock} is a generalization of BG stock models considered in \cite{bilateral2}. Recently, by incorporating a diffusion component into the BG process, Kirkby {\it et al.} \cite{bgm2024} developed a generalized exponential L\'evy model and derived an arbitrage-free option pricing formula. In practice, it is often necessary to address the appropriate pricing of European options $\Phi(S_T)$, where $T > 0$ represents the maturity time and $\Phi : \mathbb{R} \to \mathbb{R}$ denotes the payoff function. Also, the option price is given by $\pi:=e^{-rT}\mathbb{E}_{\mathbb{Q}}[\Phi(S_T)],$ where $\mathbb{Q} \sim \mathbb{P}$ is a local martingale measure. Here, $\mathbb{Q}$ is also a probability measure which is equivalent to the objective probability measure $\mathbb{P}$, such that the discounted stock price process
\begin{align}\label{dspp}
	\tilde{S}_t := e^{-(r-v)t}S_t = S_0 e^{T_n(t) - (r-v)t}, \quad t \ge 0
\end{align}
is a local $\mathbb{Q}$-martingale, see \cite{bilateral2}. 

\noindent
The following lemma follows from Lemma 3.1 of \cite{bilateral2}, which establishes necessary and sufficient conditions for the existence of a probability measure $\mathbb{Q} \sim \mathbb{P}$.
\vspace{-.5cm}
\begin{lem}
Let $r\geq v \geq 0$. Let $\{T_n(t)\}_{t\geq 0}$ be a L\'evy process with cf \eqref{cfBGRP}. Then, $\mathbb{P}$ is a martingale measure for $\tilde{S}_t$ if and only if 
\begin{align}\label{dspp1}
\mathbb{E}_{L_n} \left(\left(\frac{\eta}{\eta-1}\right)^{L_n}  \right) \mathbb{E}_{M_n} \left(\left(\frac{\xi}{\xi-1}\right)^{M_n}  \right)=(1-1/\eta)^p (1-1/\xi)^q e^{r-v},
\end{align}
where $\eta=\alpha/1-\alpha$, $\xi=\beta/1-\beta$, and $\alpha,\beta,L_n,M_n,p$, and $q$ are as defined in \eqref{ref1} and \eqref{ref2}, respectively.
\end{lem}
\vspace{-.8cm}
\begin{proof}
	From \eqref{mggd1}, we have $\mathbb{E}(e^{T_n(1)})=\mathbb{E}_{L_n} \left(\left(\frac{\eta}{\eta-1}\right)^{p+L_n}  \right) \mathbb{E}_{M_n} \left(\left(\frac{\xi}{\xi-1}\right)^{q+M_n}  \right)<\infty$. By Lemma 2.6 of \cite{bilateral2}, the discounted process $\tilde{S_t}$ in \eqref{dspp} is a local martingale if and only if $\mathbb{E}[e^{T_n(1)-(r-v)}]=1,$ which is only the case if and only if \eqref{dspp1} holds.
\end{proof}
\vspace{-.46cm}
\subsubsection{Pricing formula}
The existence of a martingale measure $\mathbb{Q} \sim \mathbb{P}$ ensures that the stock market is free of arbitrage, and the price of an European option $\Phi(S_T )$, where $T( > 0)$ is the time of maturity and $\Phi : \mathbb{R}\to \mathbb{R}$ is the pay-off function, is given by
\begin{align}
\pi=e^{-rT}\mathbb{E}_{\mathbb{Q}}[\Phi(S_T)].
\end{align}
\noindent
Also, an arbitrage-free pricing formula for a European Call Option at time $t (\ge 0)$ is, provided that $S_t = s$, given by (see \cite{bilateral0})
\begin{equation}\label{oppf1}
\pi = e^{-rT}\mathbb{E}_{\mathbb{Q}}\left[(S_T - K)^+ \mid S_t = s\right], ~T>t,
\end{equation}
where $K$ denotes the strike price. Setting $t^\prime=(T-t)$, we can express the expectation in \eqref{oppf1} as
\begin{align}\label{pf1}
\nonumber\pi&=e^{-rT}\mathbb{E}_{\mathbb{Q}}\left[(S_t e^{T_n(T)-T_n(t)} - K)^+ \mid S_t = s\right]\\
&= e^{-rT} \int_{\ln\left(\frac{K}{s}\right)}^{\infty} 
(se^x - K) h_{T_n}(x,t^\prime) dx,
\end{align}  
\noindent
where $h_{T_n}(x,t^\prime)$ denotes the pdf of $T_n(t^\prime) \sim BG \left(\eta, (L_n+p)t^\prime,\xi,(M_n+q)t^\prime \right)$ (see, Section \ref{pdfTn}). 

%Note that, from \eqref{lcbgpro1} and \eqref{denlchgf}, the pdf of $T_n(t^\prime)$ is given by
% \begin{align*}
%\nonumber & h_{T_{n}}(x,t^\prime)\\
%\nonumber&=c_nd_n\sum_{j=0}^{\infty}\gamma_j\sum_{k=0}^{\infty}\delta_k\frac{\eta^{(p+j)t^\prime} \xi^{(q+k)t^\prime}}{\Gamma (p+j)t^\prime \Gamma (q+k)t^\prime} \\
%	  &\quad \times\small{\begin{cases}
%	e^{-\eta x}x^{(p+q+j+k)t^\prime-1}\Gamma t^\prime (q+k) F((q+k)t^\prime,(p+q+j+k)t^\prime, (\eta+\xi)x),&x\in(0,\infty)\\
%	e^{\xi x}(-x)^{(p+q+j+k)t^\prime-1}\Gamma t^\prime (q+k) F((p+j)t^\prime,(p+q+j+k)t^\prime, -(\eta+\xi)x),&x\in(-\infty,0).
%	\end{cases}}
%\end{align*}
\noindent
Note that \eqref{pf1} provides a pricing formula driven by a mixed BG process, which extends the pricing formula, for the case $n\geq 2$, for BG stock models considered in Proposition 8.3 of \cite{bilateral0}.

\noindent
 As a special case, we get an arbitrage-free pricing formula driven by the L\'evy process associated with the linear combination of gamma r.v.s as:
\begin{align}\label{lcgop}
\pi
&=c_ne^{-rT}\sum_{j=0}^{\infty}\gamma_j \frac{\eta^{(p+j)t^\prime }  }{\Gamma t^\prime (p+j)}\int_{\ln\left(\frac{K}{s}\right)}^{\infty} (se^x - K)x^{(p+j)t^\prime-1}e^{-\eta x}~ dx,
\end{align}
\noindent
where $\eta=\alpha/1-\alpha$, and $c_n,\gamma_j,p, \alpha$ are defined in \eqref{ref1} and \eqref{ref2}. Let $\Gamma(w,z):=\int_{z}^{\infty}x^{w-1}e^{-x}dx,$ $~z\in \mathbb{R}$ be the incomplete gamma function. Then, \eqref{lcgop} can be written as
\begin{align*}
\pi
&=c_ne^{-rT}\sum_{j=0}^{\infty} \frac{\gamma_j   }{\Gamma  (p+j)t^\prime}\bigg( s \left(\frac{\eta}{\eta -1}\right)^{(p+j)t^\prime}  \Gamma \left((p+j)t^\prime, (\eta-1)ln\left(\frac{K}{s}\right)\right) \\
&\quad\quad-  K\Gamma \left((p+j)t^\prime, \eta ln\left(\frac{K}{s}\right)\right) \bigg).
\end{align*}
\noindent
 When $s=K$, we get an exact pricing formula as
\begin{align*}
\pi&=c_nKe^{-rT}\sum_{j=0}^{\infty}\gamma_j \frac{\eta^{(p+j)t^\prime }  }{\Gamma t^\prime (p+j)}\int_{0}^{\infty}e^xx^{(p+j)t^\prime-1}e^{-\eta x}~ dx -Ke^{-rT}\\
&=Ke^{-rT}\left(c_n\sum_{j=0}^{\infty}\gamma_j\left( \frac{\eta}{\eta -1} \right)^{(p+j)t^\prime }-1\right)\\
&=Ke^{-rT}\left( \mathbb{E}_{L_n} \left( \frac{\eta}{\eta -1} \right)^{(L_n+p)t^\prime }-1\right),
\end{align*}
\noindent
where $\mathbb{E}_{L_n}$ is the expectation with respect to the rv $L_n$.
%since $S_n(t^\prime)\sim Ga(\eta, (L_n+p)t^\prime)$.




%\noindent

%\end{rem}
 
 
 
\vskip 1ex
\noindent
\textbf{Acknowledgement:} The first author acknowledges financial support from the Research Seed Grant of NIT Warangal.
 
 
\setstretch{1}
\begin{thebibliography}{}
	
	\bibitem{apple2009} Applebaum, D. (2009). L\'evy processes and stochastic calculus, Second edition. Cambridge Studies in Advanced Mathematics, $\bf{116}$. Cambridge University Press, Cambridge, xxx+460 pp.
	
	%\bibitem{frullani} Arias-De-Reyna, J. (1990). On the Theorem of Frullani. Proceedings of the American Mathematical Society. Volume $\bf{109}$. Number I, pp. 165-175.
	
	%\bibitem{ckey0} Arteaga, A. R. and Sato, K. I (2019). Topics in Infinitely Divisible Distributions and L\'evy Processes. Springer Briefs in Probability and Mathematical Statistics, Springer.
	
	%\bibitem{Handbook1972} Abramowitz, M. and Stegun,  I.A. (1972). Handbook of Mathematical Functions. Dover Publications, New York.
	
	\bibitem{k0} Arras, B. and Houdr$\acute{e}$, C. (2019). On Stein's method for infinitely divisible laws with finite first moment. Springer Briefs in Probability and Mathematical Statistics, Springer.
	
	%\bibitem{avenaetal2025}Avena, L., Da Costa, C., \& Peterson, J. (2025). Gaussian, stable, tempered stable and mixed limit laws for random walks in cooling random environments. Annales de l'Institut Henri Poincare (B) Probabilites et statistiques. {\bf 61 (2)}, 850-889.
	
	%\bibitem{MSPN}Azmoodeh, E., and Gasbarra, D. (2019). On a new Sheffer class of polynomials related to normal product distribution. Theory of Probability and Mathematical Statistics, {\bf 98}, 51-71.
	
	
	\bibitem{MS001} Azmoodeh, E., Eichelsbacher, P. and Th$\ddot{\text{a}}$le, C. (2022). Optimal Variance-Gamma approximation on the second Wiener chaos. Journal of Functional Analysis, {\bf 282(11)}, p.109450.
	
%	\bibitem{MS0}Azmoodeh, E., Peccati, G. and Yang, X. (2021). Malliavin-Stein method: a survey of some recent developments. Modern Stochastics: Theory and Applications, \textbf{8(2)}, 141-177.
	
	%\bibitem{k22}Arras, B. and Houdr\'e, C. (2019). On Stein’s method for multivariate self-decomposable laws. Preprint
	%arXiv:1907.10050.
	
	%	\bibitem{kk4}Barbour, A. D.(1988). Stein's method and Poisson process convergence. Journal of Applied Probability $\mathbf{25A}$, pp. 175-184.
	
%	\bibitem{rosen1} Bai, S. and  Taqqu, M. (2017). Behavior of the generalized Rosenblatt process at extreme critical exponent values, Ann. Probab. \textbf{45(2)}, 1278-1324.
	
	%\bibitem{BU1} Barman, K. and Upadhye, N.S. (2024). On Stein factors for Laplace approximation and their application to random sums. Statistics \& Probability Letters, \textbf{206}, 109996.
	
	\bibitem{BUV1} Barman, K., Upadhye, N. S., \& Vellaisamy, P. (2025). Covariance Identities and Variance Bounds for Infinitely Divisible Random Variables and Their Applications. ALEA, Lat. Am. J. Probab. Math. Stat. \textbf{22}, 389-413.
	
	\bibitem{k8} Barbour, A. D. (1990). Stein's method for diffusion approximations. Probability Theory and Related Fields $\mathbf{84}$, 297-322.
	
	
\bibitem{cam2001}	Campos, L. M. B. C. (2001). On some solutions of the extended confluent hypergeometric differential equation. Journal of computational and applied mathematics, \textbf{137(1)}, 177-200.
	
	\bibitem{Ransum2024}Chadjiconstantinidis, S. (2024). Total variation distance and compound Poisson approximations for random sums. Stochastic Analysis and Applications, \textbf{42(3)}, 559-590.
	
	\bibitem{Chatterjee2007} Chatterjee, S. (2007). Lecture notes: Stein's method. Berkeley, Fall 2007.
	%\bibitem{Novak1}$\check{C}$ekanavi$\check{c}$ius, V. and Novak, S. Y. (2024). Compound Poisson approximation. CRC Press.
	
	
	%	\bibitem{k14} Boyarchenko, S.I. and Levendorskii,  S.Z. (2000). Option pricing for truncated L\'evy processes, International Journal of Theoretical and Applied Finance $\mathbf{3(3)}$, pp. 549-552.
	%	
	%	%\bibitem{privault} Breton, J.C., Privault, N. (2020). Wasserstein Distance Estimates for Stochastic Integrals by Forward-Backward Stochastic Calculus. Potential Analysis. https://doi.org/10.1007/s11118-020-09874-0
	%	
	%	\bibitem{k15}Carr, P., Geman, H., Madan,  D.B., Yor, M. (2002). The fine structure of asset returns: an empirical investigation. Journal of Business $\mathbf{75 (2)}$ pp. 305-332.
	%	
	%	
	%	
	%	\bibitem{k3}Chen, L. H. Y.(1975). Poisson approximation for dependent trials. Annals of Probability $\mathbf{3}$, pp. 534-545.
	%	
	%\bibitem{chen}Chen, L. H. Y., Goldstein, L. and Shao, Q-M. (2011). Normal Approximation by Stein's Method. Springer. 
	%	
	%	\bibitem{k16}Chen, P., Nourdin, I., Xu, L., Yang, X., Zhang, R. (2019). Non-integrable stable approximation by Stein’s method. Preprint http://arxiv.org/abs/1903.12315v2
	%	
	%	
	%	\bibitem{k17}Chen, P., Nourdin, I. and Xu, L. (2018). Stein’s method for asymmetric $\alpha-$stable distributions, with applications to CLT. Preprint https://arxiv.org/pdf/1808.02405.
	
	
		
	
	
	\bibitem{dalygaunt2016} Daly, F. and Gaunt, R.E. (2016). The Conway-Maxwell-Poisson distribution: Distributional theory and approximation. ALEA, \textbf{13}, 635-658.
	
\bibitem{DP1990} Diaconis, P. and Perlman, M. D.(1990). Bounds for tail probabilities of weighted sums of independent gamma random variables. In Topics In Statistical Dependence (IMS Lecture Notes Monogr. Ser. 16), Institute of Mathematical Statistics, Hayward, CA, pp. 147-166.	
	
%	\bibitem{HJM1999} Eberlein, E. and Raible, S. (1999). Term structure models driven by general Lévy processes. Mathematical Finance, \textbf{9(1)}, 31-53.
%	
%	\bibitem{HJM2003}Eberlein, E. and $\ddot{\text{O}}$zkan, F. (2003). The defaultable L\'evy term structure: ratings and restructuring. Mathematical Finance, 13(2), pp.277-300.
	
\bibitem{EZ1980} Engel, J. and Zijlstra, M. (1980). A characterization of the gamma distribution by the negative binomial distribution.
J. Appl. Prob. \textbf{17}, 1138-1144.	
	
	\bibitem{gamdiff}Forrester, P. J. (2024). On the gamma difference distribution. Statistics \& Probability Letters, 110136.
	
%		\bibitem{HJM1992}Heath, D., Jarrow, R. and Morton, A. (1992). Bond pricing and the term structure of interest rates: A new methodology for contingent claims valuation. Econometrica: Journal of the Econometric Society, 77-105.
%	%	
%	%	
%		\bibitem{LT1987} Jacod, J. and Shiryaev, A.N. (1987). Limit Theorems for Stochastic Processes. Springer, Berlin.
	%	
	%	
	%	\bibitem{k4}Fulman, J. and Ross, N. (2013). Exponential approximation and Stein's method of exchangeable pairs. ALEA, Latin American Journal of Probability and Mathematical Statistics $\mathbf{10(1)}$, pp. 1-13.
	%	
	%	
	%	
	%	\bibitem{N2}Gaunt, R. E. (2013). Rates of convergence of variance-gamma approximations via Stein's method. Dphil thesis, The Queen's College
	%	University of Oxford.
	
	\bibitem{Gaunt2014}Gaunt, R. E. (2014). Variance-Gamma approximation via Stein’s method. Electronic Journal of Probability 19 no. 38, 1-33.

	\bibitem{Gaunt2020}Gaunt, R.E. (2020). Wasserstein and Kolmogorov error bounds for variance-gamma approximation via Stein’s method $\textbf{I}$. Journal of Theoretical Probability $\mathbf{33}$, 465-505.
	
	\bibitem{Gaunt2022}Gaunt, R.E. (2022). Stein factors for variance-gamma approximation in the Wasserstein and Kolmogorov distances. Journal of Mathematical Analysis and Applications, \textbf{514(1)}, 126274.
	
	
	%	\bibitem{k6}Gaunt, R.E., Pickett, A.M., Reinert, G. (2017). Chi-square approximation by Stein's method with application to Pearson's statistic. Annals of Applied Probability $\mathbf{27 (2)}$, 720-756
	%	
	%	\bibitem{kk3}Gaunt, R.E. (2017). On Stein's method for products of normal random variables and zero bias couplings. Bernoulli $\mathbf{23(4B)}$, pp. 3311-3345 
	%	
	%	\bibitem{key1}Gaunt, R.E., Mijoule, G. and Swan, Y. (2019). Some new Stein operators for product distributions. Preprint https://arxiv.org/abs/1901.11460v2
	%	
	%	\bibitem{gaunt7} Gaunt, R.E. (2021). New error bounds for Laplace approximation via Stein's method. ESAIM: Probability and Statistics, \textbf{25}, pp. 325-345.
	%	%\bibitem{N1}Klebanov, L.B. and Sl\'amov\'a, L. (2015). Tempered distributions: does universal tempering procedure exist? Preprint
	%	
	%	\bibitem{kk5}Grabchak, M. (2015). Tempered stable distributions stochastic models for multiscale processes. Springer Briefs in Mathematics, Springer.
	%	
	%	
	%	\bibitem{koponen} Koponen, I. (1995). Analytic approach to the problem of convergence of truncated L\'evy flights towards the Gaussian stochastic process. Physical Review E $\mathbf{52}$. pp. 1197-1199
	%	
	%	\bibitem{k13} K$\ddot{u}$chler, U. and Tappe, S (2013). Tempered stable distributions and processes. Stochastic Stochastic Processes and their Applications $\mathbf{123}$, 4256-4293
	%	
	%	
	%	
	%	\bibitem{k19}Kumar, A.N. and Upadhye, N.S. (2017). On discrete Gibbs measure approximation to runs. Preprint https://arxiv.org/pdf/1701.03294
	%	
	%	
	%	
	%	
	%	
	%	\bibitem{k5} Luk, H. (1994). Stein's Method for the Gamma Distribution and Related Statistical Applications. PhD thesis, University of Southern California.
	\bibitem{bgm2024}Kirkby, J. L., Rinella, C. A., \& Aguilar, J. P. (2024). The bilateral Gamma motion: Calibration and option pricing. Frontiers of Mathematical Finance, \textbf{3(3)}, 400-439.
	
	\bibitem{bilateral0} K$\ddot{u}$chler, U. and Tappe, S. (2008). Bilateral gamma distributions and processes in financial mathematics. Stochastic Stochastic Processes and their Applications, $\mathbf{118(2)}$, 261-283.
	
	\bibitem{bilateral1} K$\ddot{u}$chler, U. and Tappe, S. (2008). On the shapes of bilateral Gamma densities. Statistics \& Probability Letters, \textbf{78}, 2478-2484.
	
	\bibitem{bilateral2}K$\ddot{u}$chler, U. and Tappe, S. (2009). Option pricing in bilateral Gamma stock models. Statistics \& Decisions International mathematical journal for stochastic methods and models, {\bf 27 (4)},	281-307.
	
	\bibitem{tsd2013}K$\ddot{u}$chler, U. and Tappe, S. (2013). Tempered stable distributions and processes. Stochastic processes and their applications, 123(12), 4256-4293.
	
	

		\bibitem{leyreisw2017}Ley, C., Reinert, G., Swan, Y. (2017). Stein's method for comparison of univariate distributions. Probab. Surv. \textbf{14}, pp. 1-52. 	
	
	\bibitem{nourdin} Nourdin, I. and Peccati, G. (2012). Normal approximations with Malliavin calculus. Cambridge University Press. Cambridge tracts in mathematics \textbf{192}.
	
	\bibitem{roos2003}Roos, B., and Pfeifer, D. (2003). On the distance between the distributions of random sums. Journal of applied probability, \textbf{40(1)}, 87-106.
	
	
%\bibitem{Novak1}Novak, S. Y. (2024). On Poisson Approximation. Journal of Theoretical Probability, 1-27. https://doi.org/10.1007/s10959-023-01310-4
	
	
	%	\bibitem{pike}Pike, J. and Ren, H. (2014). Stein's method and the Laplace distribution. ALEA Lat. Am. J. Probab. Math. Stat. $\bf{11}$, pp. 571–587.
	%	
	%	\bibitem{k9}Rosinski, J. (2007). Tempering stable processes. Stochastic Processes and their Applications $\mathbf{117(6)}$, pp. 677–707.
	%	
	%	\bibitem{k25}Ross, N. (2010). Fundamentals of Stein’s method. Probability Surveys $\mathbf{8}$, pp. 210-293.
	
	\bibitem{sato} Sato, K.I. (1999). L\'evy Processes and Infinitely Divisible Distributions. Cambridge University Press, Cambridge.
	
	\bibitem{sch1988} Scheuer, E. M. (1988). Reliability of an m-out-of-n system when component failure induces higher failure rates in
survivors. IEEE Trans. Reliab. {\bf 37}, 73-74.

	%	\bibitem{slepov} Slepov, N.A. (2021). Convergence rate of random geometric sum distributions to the Laplace law. Theory Probab. Appl., \textbf{66(1)}, 121-141.
	%	
	%	\bibitem{k2}Stein, C. (1972). A bound for the error in the normal approximation to the the distribution of a sum of dependent random variables. In Proceedings of the Sixth Berkeley Symposium on Mathematical Statistics and Probability, $\mathbf{2}$, Univ. California Press, Berkeley, pp. 583-602.
	%	
	%	\bibitem{den}Stein, C., Diaconis, S, Holmes, S. and Reinert, G. (2004). Use of exchangeable pairs in the analysis of simulations. In Steins method: expository lectures and applications, volume 46 of IMS Lecture Notes Monogr.Ser., Inst. Math. Statist., Beachwood, OH , pp. 1-26.
	
	\bibitem{stein} Stein, E.M. and Shakarchi, R. (2003). Fourier analysis. An introduction. Princeton Lectures in Analysis, 1. Princeton.
	
	%	\bibitem{k10}Sztonyk, P. (2010). Estimates of tempered stable densities. Journal of Theoretical Probability $\mathbf{23(1)}$, pp. 127–147.
	%	
	%	
	%	
	%	\bibitem{k23}Schoutens, W. (2001). Orthogonal polynomials in Stein’s method. Journal of Mathematical Analysis and Applications $\mathbf{253}$, pp. 515-531.
	
	
	\bibitem{PVBC} Vellaisamy, P. and Chaudhuri, B. (1996). Poisson and compound Poisson approximations for random sums of random variables. Journal of applied probability, {\bf 33(1)}, 127-137.
	
	\bibitem{pvns2009}Vellaisamy, P., \& Upadhye, N. S. (2009). On the sums of compound negative binomial and gamma random variables. Journal of Applied probability, \textbf{46(1)}, 272-283.
	
	\bibitem{k1}Upadhye, N.S. and Barman, K. (2022). A unified approach to Stein's method for stable distributions. Probability Surveys, \textbf{19}, 533-589.
	
	
	\bibitem{Villani} Villani, C. (2003). Topics in Optimal Transportation, American Mathematical Society. MR-1964483
	
	
	
	
	
	
	
	%\bibitem{k7} Upadhye, N. S., $\check{C}$ekanavi$\check{c}$ius, V. and Vellaisamy, P. (2017). On Stein operators for discrete approximations. Bernoulli $\mathbf{23}$, 2828-2859.
	
	
	
	\bibitem{VS2010} Vellaisamy, P., and Sreehari, M. (2010). Some intrinsic properties of the gamma distribution. Journal of the Japan Statistical Society, \textbf{40(1)}, 133-144.
	
	
	
	
	
	
	
	
	
	
	%\bibitem{k12}Uwe K$\ddot{u}$chler, Stefan Tappe(2008), Bilateral gamma distributions and processes in financial mathematics, Stochastic Processes and their Applications 118, 261–283
	
	
	%	\bibitem{k20}Xu, L. (2019). Approximation of stable law in Wasserstein-1 distance by Stein’s method. The Annals of Applied Probability $\mathbf{29}$, No. 1, pp. 458-504.
		%\bibitem{sato1999} Sato, K.I. (1999). L\'evy Processes and Infinitely Divisible Distributions. Cambridge University Press, Cambridge.
\end{thebibliography}



\end{document}











\section*{Discounted Stock Price as a Local $\mathbb{Q}$-Martingale \cite{bilateral2}}

\begin{defn}
	Let $(\Omega,\mathcal{F},\mathbb{Q})$ be a probability space and 
	$(X_t)_{t\ge 0}$ a stochastic process such that 
	$\mathbb{E}_{\mathbb{Q}}[|X_t|] < \infty$ for all $t \ge 0$.
	The process $(X_t)_{t\ge 0}$ is called a local $\mathbb{Q}$-martingale if, 
	for all $0 \le s \le t$,
	\[
	\mathbb{E}_{\mathbb{Q}}\!\left[X_t \mid X_u,\; 0 \le u \le s \right] = X_s 
	\quad \text{a.s.}
	\]
\end{defn}
\noindent
Some important lemmas:
\begin{lem}
	Let $\{X_t\}_{t\geq 0}$ be a real-valued L\'evy process. The following statements are equivalent:
	\begin{enumerate}
		\item $X$ is a local martingale;
		\item $X$ is a martingale;
		\item $\mathbb{E}[X_1] = 0$.
	\end{enumerate}
\end{lem}

\begin{lem}
	Let $\{X_t\}_{t\geq 0}$ be a real-valued L\'evy process. The following statements are equivalent:
	\begin{enumerate}
		\item $e^{X}$ is a local martingale;
		\item $e^{X}$ is a martingale;
		\item $\mathbb{E}\!\left[e^{X_1}\right] = 1$.
	\end{enumerate}
\end{lem}

 An exponential L\'evy stock model is of the form
$$
S_t = S_0 e^{X_t}, \qquad B_t = e^{rt},
$$
where
\begin{itemize}
	\item $S_t$ is the stock price,
	\item $S_0 > 0$ is the initial stock price,
	\item $X_t$ is a Lévy process,
	\item $B_t$ is the bank account,
	\item $r \ge 0$ is the constant interest rate,
	\item $q \ge 0$ is the continuous dividend rate (investors receive part of return as dividends).
\end{itemize}

\subsection*{Discounted Stock Price}

Since the stock pays continuous dividends at rate $q$, the discounted stock price process is
\[
\tilde{S}_t = e^{-rt} S_t e^{qt}.
\]
Equivalently,
\[
\tilde{S}_t = e^{-(r-q)t} S_t.
\]

Substituting $S_t = S_0 e^{X_t}$ gives
\[
\tilde{S}_t = S_0 e^{X_t - (r-q)t}.
\]

\subsection*{Risk-Neutral Martingale Condition}

Under an equivalent probability measure $\mathbb{Q} \sim \mathbb{P}$, the absence of arbitrage requires that the discounted stock price process be a local martingale. 

This implies
\[
\mathbb{E}_{\mathbb{Q}} \left[ e^{X_t - (r-q)t} \right] = 1,
\]
or equivalently,
\[
\mathbb{E}_{\mathbb{Q}} \left[ e^{X_t} \right] = e^{(r-q)t}.
\]

Setting $t=1$, we obtain the martingale condition
\[
\mathbb{E}_{\mathbb{Q}} \left[ e^{X_1} \right] = e^{r-q}.
\]

\subsection*{Lévy Exponent Formulation}

If $X_t$ has Lévy exponent $\psi(u)$ defined by
\[
\mathbb{E}\left[e^{uX_t}\right] = e^{t\psi(u)},
\]
then the martingale condition becomes
\[
e^{t\psi(1)} = e^{(r-q)t}.
\]
Hence,
\[
\psi(1) = r - q.
\]

\subsection*{Interpretation}

Under the risk-neutral probability measure $\mathbb{Q}$:
\begin{itemize}
	\item The drift of the Lévy process is adjusted,
	\item The expected growth rate of the stock becomes $r-q$,
	\item The discounted stock price process
	\[
	\tilde{S}_t = S_0 e^{X_t - (r-q)t}
	\]
	is a local $\mathbb{Q}$-martingale.
\end{itemize}




\noindent
\section*{Bilateral Gamma Process (K$\ddot{u}$chler and Tappe \cite{bilateral0})} Recall that bilateral gamma (BG) distributions are infinitely divisible. Let us list the properties of the associated L\'evy processes, which are denoted by $X$ in the sequel. From the representation (2.3) of the L\'evy measure $\nu_{bg}$ we see that $\nu_{bg}(\mathbb{R}) = \infty$ and $\int_{\mathbb{R}} |x| \nu_{bg}(dx) < \infty.$ We now have the following properties. 
 
\begin{enumerate}
\item BG processes are finite-variation processes, making infinitely many jumps at each interval with positive length, and they are equal to the sum of their jumps, i.e.
\[
X_t = \sum_{s \leq t} \Delta X_s = x * \mu^X, \quad t \geq 0
\]
where $\mu^X$ denotes the random measure of jumps of $X$.

\item BG processes are special semimartingales with canonical decomposition
\[
X_t = x * (\mu^X - \zeta)_t + \left( \frac{\alpha^+}{n^+} - \frac{\alpha^-}{n^-} \right) t, \quad t \geq 0
\]
where $\zeta$ is the compensator of $\mu^X$, which is given by $\zeta(dt, dx) = dt \nu_{bg}(dx)$ with $\nu_{bg}$ denoting the L\'evy measure given by (2.3).

\item We immediately see from the characteristic function (2.2) that all increments of $X$ have a bilateral Gamma distribution, or more precisely
$$X_t - X_s \sim BG(\alpha, p(t-s), \beta, q(t-s)), \quad \text{for } 0 \leq s < t. $$

\end{enumerate} 
 
 \noindent
We assume that the probability space $(\Omega, \mathcal{F}, \mathbb{P})$ is given as follows. 
Let $\Omega = \mathbb{D}$, the collection of functions $\omega(t)$ from $\mathbb{R}_{+}$ into $\mathbb{R}$, 
right-continuous with left limits. For $\omega \in \Omega$, let $X_t(\omega) = \omega(t)$ and let 
$\mathcal{F} = \sigma(X_t : t \in \mathbb{R}_{+})$ and $(\mathcal{F}_t)_{t \geq 0}$ be the filtration 
$\mathcal{F}_t = \sigma(X_s : s \in [0,t])$. We consider a probability measure $\mathbb{P}$ on $(\Omega, \mathcal{F})$ 
such that $X$ is a BG process.

\subsection{Arbitrage-free Stock models}
We move on to present some applications to finance of the theory developed above. 
Assume that the evolution of an asset price is described by an exponential L\'evy model 
\[
S_t = S_0 e^{rt + X_t},
\]
where $S_0 > 0$ is the (deterministic) initial value of the stock, $r$ the interest rate, 
and where the L\'evy process $X$ is a BG process under the measure $\mathbb{P}$, 
which plays the role of the real-world measure. 

\vskip 2ex
\noindent
In order to avoid arbitrage, the question of whether there exists an equivalent martingale measure arises, 
i.e. is there a measure $\mathbb{Q} \overset{^{\text{loc}}}{ \sim} \mathbb{P}$ such that 
\[
Y_t := e^{-rt} S_t
\]
is a local martingale?


\begin{lem}
Assume $\alpha > 1$. Then $Y$ is a local $\mathbb{P}$-martingale if and only if
\begin{align}\label{bglmar2}
\left( \frac{\alpha}{\alpha-1} \right)^{p}
= \left( \frac{\beta+1}{\beta} \right)^{q} .
\end{align}
\end{lem}

\begin{proof}
 Since the Gaussian part of the BG process $X$ is zero, 
It\^o’s formula \cite[Thm. I.4.57]{LT1987}, applied on $Y_t = S_0 e^{X_t}$, yields for the discounted stock prices
\[
Y_t = Y_0 + \int_0^t Y_{s-} dX_s + S_0 \sum_{0 < s \leq t} 
\left( e^{X_s} - e^{X_{s-}} - e^{X_{s-}} \Delta X_s \right).
\]

\noindent
Recall that $X = x * \mu^X$ and that the compensator $\zeta$ of $\mu^X$ is given by 
$\zeta(dt, dx) = dt\,\nu_{bg}(dx)$, where $\nu_{bg}$ denotes the L\'evy measure from (2.3). So we obtain
\begin{align}
\nonumber Y_t &= Y_0 + \int_0^t \int_{\mathbb{R}} x Y_{s-} \mu^X(ds, dx) 
     + \int_0^t \int_{\mathbb{R}} Y_{s-}(e^x - 1 - x)\mu^X(ds, dx) \\
&= Y_0 + \int_0^t \int_{\mathbb{R}} Y_{s-}(e^x - 1)(\mu^X - \zeta)(ds, dx) 
     + \int_0^t Y_s \int_{\mathbb{R}} (e^x - 1) \nu_{bg}(dx)\, ds. \label{bglmar1}
\end{align}

\noindent
Applying Frullani improper integral formula, the integral in the drift term is equal to
\begin{align}
\int_{\mathbb{R}} (e^x - 1) \nu_{bg}(dx) 
\nonumber&= p \int_0^\infty \frac{e^{-(\alpha - 1)x} - e^{-\alpha x}}{x}\, dx 
- q \int_0^\infty \frac{e^{-\beta x} - e^{-(\beta+1)x}}{x}\, dx\\
&= p \ln \left( \frac{\alpha}{\alpha - 1} \right) 
- q\ln \left( \frac{\beta + 1}{\beta} \right). \label{bglmar}
\end{align}

\noindent
The discounted price process $Y$ is a local martingale if and only if the drift in \eqref{bglmar1} vanishes, 
and by \eqref{bglmar} this is the case if and only if \eqref{bglmar2} is satisfied. 
\end{proof}

\noindent
As is usual in financial modelling with jump processes, the market is free of arbitrage, but not complete; that is, there exist several martingale measures. The next result shows that we can find a continuum 
of martingale measures by staying within the class of BG processes. We define $\phi : (1,\infty) \to \mathbb{R}$ as


\[
\phi(\alpha) := \phi(\alpha; p, q) := 
\left( \left( \frac{\alpha}{\alpha-1} \right)^{p / q} - 1 \right)^{-1}, 
\quad \alpha \in (1,\infty).
\]

\begin{pro}
For each $\alpha \in (1,\infty)$, there exists a martingale measure 
$\mathbb{Q}_\alpha \overset{^{\text{loc}}} \sim \mathbb{P}$ such that under $\mathbb{Q}_\alpha$ we have
\begin{align}\label{eqmeas}
X_1 \sim BG(\alpha, p , \phi(\alpha), q). 
\end{align}
\end{pro}

\begin{proof}
 Recall that $X$ is $BG(\alpha,p,\beta,q)$ under $\mathbb{P}$. According to Proposition 6.1 of \cite{bilateral0}, for each $\alpha \in (1,\infty)$ there exists a probability measure $\mathbb{Q}_\alpha \overset{{\text{loc}}}{ \sim} \mathbb{P}$ such that under $\mathbb{Q}_\alpha$ relation \eqref{eqmeas} is fulfilled, and moreover this measure $\mathbb{Q}_\alpha$ is a martingale measure, because Euation \eqref{bglmar2} from Lemma 5.2 is satisfied. 

\end{proof}

%\textbf{Proof.}




\subsection{Term structure models}
Let $f(t, T)$ be a Heath-Jarrow-Morton (HJM) term structure model \cite{HJM1992}
$$
\mathrm{d} f(t, T)=\alpha(t, T) \mathrm{d} t+\sigma(t, T) \mathrm{d} X_t
$$
\noindent
driven by a one-dimensional Lévy process $X_t$. We assume that the cumulant generating function $\psi$ exists on some non-void closed interval $I \subset \mathbb{R}$ having zero as inner point. By Equation (2.1), this condition is satisfied for BG processes, by taking any non-void closed interval $I \subset\left(-\beta, \alpha\right)$ with zero as an inner point. Recall that, the cumulating generating function is given by
$$\psi(z)= \log\bigg(\bigg(\frac{\alpha}{\alpha-z}\bigg)^{p} \bigg(\frac{\beta}{\beta+z} \bigg)^{q}\bigg),~z\in\mathbb{R}.$$
\noindent
 We assume that the volatility $\sigma$ is deterministic and that, in order to avoid arbitrage, the drift $\alpha$ satisfies the HJM drift condition



$$
\alpha(t, T)=-\sigma(t, T) \psi^{\prime}(\Sigma(t, T)), \quad \text { where } \Sigma(t, T)=-\int_t^T \sigma(t, s) \mathrm{ds}
$$
\noindent
This condition on the drift is, for instance, derived in \cite[ Sec. 2.1]{HJM2003}. Since $\psi$ is only defined on $I$, we impose the additional condition
\begin{align}\label{tsm1}
\Sigma(t, T) \in I \quad \text { for all } 0 \leq t \leq T.
\end{align}
\noindent
It was shown in \cite{HJM1999} that the short rate process $r_t=f(t, t)$ is a Markov process if and only if the volatility factorizes, i.e. $\sigma(t, T)=\tau(t) \zeta(T)$. Moreover, provided the differentiability of $\tau$ as well as $\tau(t) \neq 0, t \geq 0$ and $\zeta(T) \neq 0, T \geq 0$, there exists an affine one-dimensional realization. Since $\sigma(\cdot, T)$ satisfies for each fixed $T \geq 0$ the ordinary differential equation
$$
\frac{\partial}{\partial t} \sigma(t, T)=\frac{\tau^{\prime}(t)}{\tau(t)} \sigma(t, T), \quad t \in[0, T].
$$
\noindent
We verify by using It$\hat{\text{o}}$'s formula \cite[Thm. I.4.57]{LT1987} for fixed $T \geq 0$ that such a realization
\begin{align}\label{tsm2}
f(t, T)=a(t, T)+b(t, T) Z_t, \quad 0 \leq t \leq T
\end{align}
\noindent
is given by
\begin{align}\label{tsm3}
a(t, T)=f(0, T)+\int_0^t \alpha(s, T) \mathrm{d} s, \quad b(t, T)=\sigma(t, T)
\end{align}
\noindent
and the one-dimensional state process $Z_t$, which is the unique solution of the stochastic differential equation
$$
\left\{\begin{array}{l}
d Z_t=-\frac{\tau^{\prime}(t)}{\tau(t)} Z_t d t+\mathrm{d} X_t \\
Z_0=0
\end{array}\right.
$$

\noindent
We can transform this realization into an affine short rate realization. By \eqref{tsm2}, it holds for the short rate $r_t=a(t, t)+b(t, t) Z_t, t \geq 0$, implying

$$
Z_t=\frac{r_t-a(t, t)}{b(t, t)}, \quad t \geq 0.
$$
\noindent
Inserting this equation into \eqref{tsm2}, we get

$$
f(t, T)=a(t, T)+\frac{b(t, T)}{b(t, t)}\left(r_t-a(t, t)\right), \quad 0 \leq t \leq T
$$

\noindent
Incorporating \eqref{tsm3}, we arrive at

\begin{align}\label{tsm4}
\nonumber f(t, T)= & f(0, T)-\int_0^t\left[\psi^{\prime}(\Sigma(s, T))-\psi^{\prime}(\Sigma(s, t))\right] \sigma(s, T) d s \\
& +\frac{\zeta(T)}{\zeta(t)}\left(r_t-f(0, t)\right).
\end{align}

\noindent
As an example, let $f(t, T)$ be a term structure model having a Vasi$\hat{\text{c}}$ek volatility structure, i.e.
\begin{align}\label{tsm5}
\sigma(t, T)=-\hat{\sigma} \mathrm{e}^{-a(T-t)}, \quad 0 \leq t \leq T.
\end{align}
\noindent
with real constants $\hat{\sigma}>0$ and $a \neq 0$. We assume that $a>0$ and $\frac{\hat{\sigma}}{a}<\frac{\alpha}{1-\alpha}$. Since

\begin{align}\label{tsm6}
\Sigma(t, T)=\frac{\hat{\sigma}}{a}\left(1-e^{-a(T-t)}\right), \quad 0 \leq t \leq T.
\end{align}
\noindent
We find a suitable interval $I \subset\left(-\beta, \alpha \right)$ such that condition \eqref{tsm1} is satisfied. By the results above, the short rate $r$ is a Markov process and there exists a short rate realization. Equation \eqref{tsm4} simplifies to

\begin{align}\label{tsm7}
\nonumber f(t, T)= & f(0, T)+\psi(\Sigma(0, T))-\psi(\Sigma(t, T))-e^{-a(T-t)} \psi(\Sigma(0, t)) \\
& +e^{-a(T-t)}\left(r_t-f(0, t)\right).
\end{align}

\noindent
We can compute the bond prices $P(t, T)$ by using the following result.
\begin{pro}
It holds for the bond prices
$$
P(t, T)=e^{\phi_1(t, T)-\phi_2(t, T) r_t}, \quad 0 \leq t \leq T
$$
where the functions $\phi_1, \phi_2$ are given by



\begin{align}
\nonumber\phi_1(t, T)= & -\int_t^T f(0, s) \mathrm{d} s-\int_t^T \psi\left(\frac{\hat{\sigma}}{a}\left(1-\mathrm{e}^{-a s}\right)\right) \mathrm{d} s \\
\nonumber& +\int_t^T \psi\left(\frac{\hat{\sigma}}{a}\left(1-\mathrm{e}^{-a(s-t)}\right)\right) \mathrm{d} s \\
\label{tsm8}& +\frac{1}{a}\left(1-\mathrm{e}^{-a(T-t)}\right)\left[f(0, t)+\psi\left(\frac{\hat{\sigma}}{a}\left(1-\mathrm{e}^{-a t}\right)\right)\right], \\
\label{tsm9}\phi_2(t, T)= & \frac{1}{a}\left(1-\mathrm{e}^{-a(T-t)}\right).
\end{align}	
\end{pro} 
\begin{proof}
	The claimed formula for the bond prices follows from the identity
	$$
	P(t, T)=\mathrm{e}^{-\int_1^T f(t, s) d t}
	$$
	and Equations \eqref{tsm6} and \eqref{tsm7}.
\end{proof} 




\subsection{Term structure models}
Let $f(t, T)$ be a Heath-Jarrow-Morton (HJM) term structure model \cite{HJM1992}
$$
\mathrm{d} f(t, T)=\alpha(t, T) \mathrm{d} t+\sigma(t, T) \mathrm{d} X_t
$$
\noindent
driven by a one-dimensional Lévy process $X_t$. We assume that the cumulant generating function $\psi$ exists on some non-void closed interval $I \subset \mathbb{R}$ having zero as inner point. By Equation \eqref{lcbgcgf}, this condition is satisfied for BG processes, but with random parameters, by taking any non-void closed interval $I \subset\left(-\frac{\beta}{1-\beta}, \frac{\alpha}{1-\alpha}\right)$ with zero as an inner point. Recall that, the cumulating generating function (see, \eqref{lcbgcgf}) is given by
$$\psi(z)= \log\bigg(\mathbb{E}_1\bigg(\frac{\eta}{\eta-z}\bigg)^{p+L_n} \mathbb{E}_2\bigg(\frac{\xi}{\xi+z} \bigg)^{q+M_n}\bigg),~z\in\mathbb{R},$$
\noindent
where $\eta=\frac{\alpha}{1-\alpha},$ and $\xi=\frac{\beta}{1-\beta}$. We assume that the volatility $\sigma$ is deterministic and that, in order to avoid arbitrage, the drift $\alpha$ satisfies the HJM drift condition



$$
\alpha(t, T)=-\sigma(t, T) \psi^{\prime}(\Sigma(t, T)), \quad \text { where } \Sigma(t, T)=-\int_t^T \sigma(t, s) \mathrm{ds}
$$
\noindent
This condition on the drift is, for instance, derived in \cite[ Sec. 2.1]{HJM2003}. Since $\psi$ is only defined on $I$, we impose the additional condition
\begin{align}\label{LCtsm1}
\Sigma(t, T) \in I \quad \text { for all } 0 \leq t \leq T.
\end{align}
\noindent
It was shown in \cite{HJM1999} that the short rate process $r_t=f(t, t)$ is a Markov process if and only if the volatility factorizes, i.e. $\sigma(t, T)=\tau(t) \zeta(T)$. Moreover, provided the differentiability of $\tau$ as well as $\tau(t) \neq 0, t \geq 0$ and $\zeta(T) \neq 0, T \geq 0$, there exists an affine one-dimensional realization. Since $\sigma(\cdot, T)$ satisfies for each fixed $T \geq 0$ the ordinary differential equation
$$
\frac{\partial}{\partial t} \sigma(t, T)=\frac{\tau^{\prime}(t)}{\tau(t)} \sigma(t, T), \quad t \in[0, T].
$$
\noindent
We verify by using It$\hat{\text{o}}$'s formula \cite[Thm. I.4.57]{LT1987} for fixed $T \geq 0$ that such a realization
\begin{align}\label{LCtsm2}
f(t, T)=a(t, T)+b(t, T) Z_t, \quad 0 \leq t \leq T
\end{align}
\noindent
is given by
\begin{align}\label{LCtsm3}
a(t, T)=f(0, T)+\int_0^t \alpha(s, T) \mathrm{d} s, \quad b(t, T)=\sigma(t, T)
\end{align}

\noindent
and the one-dimensional state process $Z_t$, which is the unique solution of the stochastic differential equation
$$
\left\{\begin{array}{l}
d Z_t=-\frac{\tau^{\prime}(t)}{\tau(t)} Z_t d t+\mathrm{d} X_t \\
Z_0=0
\end{array}\right.
$$


\noindent
We can transform this realization into an affine short rate realization. By \eqref{LCtsm2}, it holds for the short rate $r_t=a(t, t)+b(t, t) Z_t, t \geq 0$, implying

$$
Z_t=\frac{r_t-a(t, t)}{b(t, t)}, \quad t \geq 0.
$$
\noindent
Inserting this equation into \eqref{LCtsm2}, we get

$$
f(t, T)=a(t, T)+\frac{b(t, T)}{b(t, t)}\left(r_t-a(t, t)\right), \quad 0 \leq t \leq T
$$

\noindent
Incorporating \eqref{LCtsm3}, we arrive at

\begin{align}\label{LCtsm4}
\nonumber f(t, T)= & f(0, T)-\int_0^t\left[\psi^{\prime}(\Sigma(s, T))-\psi^{\prime}(\Sigma(s, t))\right] \sigma(s, T) d s \\
& +\frac{\zeta(T)}{\zeta(t)}\left(r_t-f(0, t)\right).
\end{align}

\noindent
As an example, let $f(t, T)$ be a term structure model having a Vasi$\hat{\text{c}}$ek volatility structure, i.e.
\begin{align}\label{LCtsm5}
\sigma(t, T)=-\hat{\sigma} \mathrm{e}^{-a(T-t)}, \quad 0 \leq t \leq T.
\end{align}
\noindent
with real constants $\hat{\sigma}>0$ and $a \neq 0$. We assume that $a>0$ and $\frac{\hat{\sigma}}{a}<\frac{\alpha}{1-\alpha}$. Since

\begin{align}\label{LCtsm6}
\Sigma(t, T)=\frac{\hat{\sigma}}{a}\left(1-e^{-a(T-t)}\right), \quad 0 \leq t \leq T.
\end{align}
\noindent
We find a suitable interval $I \subset\left(-\frac{\beta}{1-\beta}, \frac{\alpha}{1-\alpha}\right)$ such that condition \eqref{LCtsm1} is satisfied. By the results above, the short rate $r$ is a Markov process and there exists a short rate realization. Equation \eqref{LCtsm4} simplifies to

\begin{align}\label{LCtsm7}
\nonumber f(t, T)= & f(0, T)+\psi(\Sigma(0, T))-\psi(\Sigma(t, T))-e^{-a(T-t)} \psi(\Sigma(0, t)) \\
& +e^{-a(T-t)}\left(r_t-f(0, t)\right).
\end{align}

\noindent
We can compute the bond prices $P(t, T)$ by using the following result.
\begin{pro}
It holds for the bond prices
$$
P(t, T)=e^{\phi_1(t, T)-\phi_2(t, T) r_t}, \quad 0 \leq t \leq T
$$
where the functions $\phi_1, \phi_2$ are given by
\begin{align}
\nonumber\phi_1(t, T)= & -\int_t^T f(0, s) \mathrm{d} s-\int_t^T \psi\left(\frac{\hat{\sigma}}{a}\left(1-\mathrm{e}^{-a s}\right)\right) \mathrm{d} s \\
\nonumber& +\int_t^T \psi\left(\frac{\hat{\sigma}}{a}\left(1-\mathrm{e}^{-a(s-t)}\right)\right) \mathrm{d} s \\
\label{LCtsm8}& +\frac{1}{a}\left(1-\mathrm{e}^{-a(T-t)}\right)\left[f(0, t)+\psi\left(\frac{\hat{\sigma}}{a}\left(1-\mathrm{e}^{-a t}\right)\right)\right], \\
\label{LCtsm9}\phi_2(t, T)= & \frac{1}{a}\left(1-\mathrm{e}^{-a(T-t)}\right).
\end{align}	
\end{pro} 
\begin{proof}
	The claimed formula for the bond prices follows from the identity
	$$
	P(t, T)=\mathrm{e}^{-\int_1^T f(t, s) d t}
	$$
	and Equations \eqref{LCtsm6} and \eqref{LCtsm7}.
\end{proof} 


\noindent
The next result gives an error bound between two compound Poisson distributions. A similar bound can be found for the total variation distance in \cite{roos2003}.

\begin{lem}\label{TVDis}
Let $Z=\sum_{i=1}^{N}Z_i$, where  $(Z_i)_{i\geq 1}$ is a sequence of iid r.v.s with cf $\phi_{T_m}^{\frac{1}{n}}(z)$ and $N\sim Poi (\lambda)$ which is independent of $Z_i$. Then, for any $j\in \mathbb{Z}_+$,
\begin{align}\label{tv0}
d_K(\Tilde{Z},Z) \leq \sum_{k=0}^{\infty}|j-k| |\mathbb{P}(M=k) - \mathbb{P}(N=k) |.
\end{align}
\end{lem}
\begin{proof}
Under the given assumption, from Theorem 1 of \cite{roos2003}, we have
\begin{align}\label{tv1}
d_{TV}(\Tilde{Z},Z) \leq \sum_{k=0}^{\infty}|j-k| |\mathbb{P}(M=k) - \mathbb{P}(N=k) |.
\end{align}
\noindent
Also, by \eqref{reltvk}, we have
\begin{align}\label{tv2}
d_{K}(\Tilde{Z},Z)  \leq d_{TV}(\Tilde{Z},Z) .
\end{align}
\noindent
Using \eqref{tv1} in \eqref{tv2}, the bound \eqref{tv0} holds.
\end{proof}
\noindent
Using the above two lemmas, we can now deduce, for the Kolmogorov distance $d_K$, an error bound for compound Poisson approximation of linear combination of BG r.v.s. 

\begin{thm}
For $w_j^{(1)},w_j^{(2)}>0$ and $i\in \mathbb{Z}_+ \setminus \{ 0\}$, let $T_m= \sum_{i=1}^{m}(w_j^{(1)}X_j-w_j^{(2)}Y_j)$ be defined in Theorem \ref{THM1}. Then, for any $j\in \mathbb{Z}_+$,


	\begin{align}\label{finalb}
	\nonumber d_{K}(T_m,Z) \leq &~ c \bigg(\sum_{j=1}^{2}|C_{j}(T_m)|\bigg)^{\frac{2}{5}} \bigg(\frac{1}{n}\bigg)^{\frac{1}{5}} \\
	&\quad \quad + \sum_{k=0}^{\infty}|j-k| |\mathbb{P}(M=k) - \mathbb{P}(N=k) |,
	\end{align}
	where $c>0$ is some positive constant and $M,N$ are r.v.s as defined in Lemma \ref{TVDis}. 
\end{thm}
\begin{proof}
By triangle inequality, we have
\begin{align}\label{DKtri}
d_K(T_m,Z) \leq d_K(T_m,\Tilde{Z}) +d_K(\Tilde{Z},Z).
\end{align}
\noindent
Using \eqref{Kbound1} and \eqref{tv0} in \eqref{DKtri}, we get \eqref{finalb}, as desired.
\end{proof}


\end{document}








